\documentclass[11pt,a4paper]{article}
\pdfoutput=1
\usepackage{jheppub}
\usepackage{amsmath}
%\input epsf
%\usepackage{epsfig}
\usepackage{amssymb}
\usepackage{graphics}
\usepackage[active]{srcltx}
%\usepackage{epstopdf}
\usepackage{pdfsync}
\usepackage{shuffle}
\usepackage{slashed}

\usepackage{subfigure}
%\usepackage{pdfsync}

\setlength{\oddsidemargin}{0.75in}
\setlength{\evensidemargin}{0.75in} \setlength{\topmargin}{0.75in}
\setlength{\textwidth}{7.0in} \setlength{\textheight}{8.5in}
\renewcommand{\baselinestretch}{1.2}
\jot=2mm


%\newcommand{\be}{\begin{equation}}
%\newcommand{\ee}{\end{equation}}
\newcommand{\bea}{\begin{eqnarray}}
\newcommand{\eea}{\end{eqnarray}}
%\newcommand{\bean}{\begin{eqnarray*}}
%\newcommand{\eean}{\end{eqnarray*}}
\newcommand{\nn}{\nonumber \\}
%\newcommand{\mat}[1]{\left( \matrix{#1} \right)}
%\newcommand{\tmat}[1]{{\scriptsize \mat{#1}}}
%\newtheorem{theorem}{\sf THEOREM}
%\newcommand{\bbibitem}[1]{\bibitem{#1}\marginpar{#1}}
%\def\thetheorem{\thesection.\arabic{theorem}}
\DeclareMathOperator{\pslash}{\displaystyle{\not}}


\def\IZ{\mathbb{Z}}
\def\IR{\mathbb{R}}
\def\IQ{\mathbb{Q}}
\def\IC{\mathbb{C}}
\def\IP{\mathbb{P}}
\def\O #1{\overline{#1}}
\def\D #1{\dot{#1}}
\def\W #1{\widetilde{#1}}
\def\WH #1{\widehat{#1}}
\def\und #1{\underline{#1}}

\def\rightaction#1{\stackrel{\rightarrow}{\partial \over \partial #1}}
\def\leftaction#1{\stackrel{\lefttarrow}{\partial \over \partial #1}}
\def\func#1{\mathop{\rm #1}\nolimits}
\def\abs#1{\left| #1\right|}
\def\braket#1{\left\langle #1 \right\rangle}
\def\bra#1{\left\langle #1\right|}
\def\ket#1{\left| #1\right\rangle}
\def\gb #1{ \left\langle #1 \right]}
\def\tgb #1{ \left[ #1 \right\rangle}


\def\bbar#1{ \overline #1}
\def\Tr{\mathop{\rm Tr}}
\def\det{\mathop{\rm det}}
\newcommand{\fref}[1]{Figure~\ref{#1}}
\def\eref#1{(\ref{#1})}
\def\bit#1{\begin{array}{c} \\ \\ \end{array} \hspace{#1 cm}}
\def\d{{\rm d}}
\def\wt{\widetilde}
\def\wtl{\widetilde{\lambda}}
\def\wh{\widehat}
\def\wht{\widehat{x}^0}
\def\whx{\widehat{x}^1}
\def\th{{\theta}}
\def\bth{{\overline{\theta}}}
\def\a{{\alpha}}
\def\ba{{\overline{\alpha}}}
\def\da{{\dot{\alpha}}}
\def\b{{\beta}}
\def\db{{\dot{\beta}}}
\def\c{{\gamma}}
\def\dc{{\dot{\gamma}}}

\def\d{\partial}
\def\rmd{{\rm d}}
\def\la{\lambda}
\def\eps{\epsilon}
\def\lblb{({\bar\la}{\bar\la})}
\def\ald{{\dot\alpha}}
\def\bed{{\dot\beta}}
\def\gad{{\dot\gamma}}
\def\sid{{\dot\rho}}
\def\vev{\braket}
\def\tgb #1{ \left[ #1 \right\rangle}
\def\bra#1{\left\langle #1\right|}
\def\bbra#1{\left[ #1\right|}
\def\ket#1{\left| #1\right\rangle}
\def\bket#1{\left| #1\right]}
\def\bvev#1{\left[ #1 \right]}
\def\Spaa{\vev}
\def\Spbb{\bvev}
\def\Spab{\gb}
\def\Spba{\tgb}




%%%%%%%%%%%%%%%%%%%%%%%%%%%%%%%%%


\title{Expanding single trace YMS amplitudes with gauge invariant coefficients}
\author[a]{Fang-Stars Wei,}
\author[b]{Kang Zhou,}

\affiliation[a]{Center for Gravitation and Cosmology, College of Physical Science and Technology, Yangzhou University,\\
No.180, Siwangting Road, Yangzhou, 225009, P.R. China.}
\affiliation[b]{Center for Gravitation and Cosmology, College of Physical Science and Technology, Yangzhou University,\\
No.180, Siwangting Road, Yangzhou, 225009, P.R. China.}


\emailAdd{mx120220339@stu.yzu.edu.cn}
\emailAdd{zhoukang@yzu.edu.cn}


\date{\today}
\abstract{In this note, we use the soft bootstrap method to obtain the expansion of single trace Yang-Mills-scalar amplitudes which manifests the gauge invariance for any polarization carried by external gluons. Our result is equivalent to the expansion proposed by Clifford Cheung and James Mangan via the so called covariant double copy framework.
}



%\keywords{Soft theorem, Soft factor, Amplitude, Expansion}

\begin{document}


\maketitle \flushbottom

%%%%%%%%%%%%%%%%%%%%%%%%%%%%%%%%%%%%%%
\section{Introduction}
\label{sec-intro}
%%%%%%%%%%%%%%%%%%%%%%%%%%%%%%

The investigations of scattering amplitudes in the past decade revealed that the tree amplitudes can be linked via the expansions \cite{Stieberger:2016lng,Schlotterer:2016cxa,Chiodaroli:2017ngp,DelDuca:1999rs,Nandan:2016pya,delaCruz:2016gnm,Fu:2017uzt,Teng:2017tbo,
Du:2017kpo,Du:2017gnh,Feng:2019tvb,Zhou:2022orv}.
For instance, the single trace Yang-Mills-scalar (YMS) amplitude ${\cal A}_{\rm YS}(1,\cdots,n;p_1,\cdots,p_m|\sigma_{n+m})$, which includes $n$ external scalars labeled by $i\in\{1,\cdots,n\}$, and $m$
external gluons encoded by $j\in\{p_1,\cdots,p_m\}$, can be expanded to YMS ones with less external gluons and more scalars,
%
\bea
& &{\cal A}_{\rm YS}(1,\cdots,n;p_1,\cdots,p_m|\sigma_{n+m})\nn
&=&\sum_{\vec{\pmb{\a}}}\,(\epsilon_{p_1}\cdot F_{\vec{\pmb{\a}}}\cdot Y_{\vec{\pmb{\a}}})\,{\cal A}_{\rm YS}(1,\{2,\cdots,n-1\}\shuffle\{\vec{\pmb{\a}},p_1\},n;\{p_2,\cdots,p_m\}\setminus{\pmb{\a}}|\sigma_{n+m})\,.~~\label{expan-YMS}
\eea
%
Here the YMS amplitude ${\cal A}_{\rm YS}(1,\cdots,n;p_1,\cdots,p_m|\sigma_{n+m})$ carries two color orderings, one is $\sigma_{n+m}$ among all external legs, another one is $(1,\cdots,n)$ among only external scalars. Two orderings are separated by $|$, while scalars and gluons are separated by $;$. At the second line of \eref{expan-YMS}, $\pmb{\a}$ are subsets of external legs $\{p_2,\cdots,p_m\}$ and are allowed to be empty, and the ordered subsets $\vec{\pmb{\a}}$ are obtained by giving orders to elements in $\pmb{\a}$. Both two sets $\{2,\cdots,n-1\}$ and $\{\vec{\pmb{\a}},p_1\}$ in $\{2,\cdots,n-1\}\shuffle\{\vec{\pmb{\a}},p_1\}$ are ordered, where $\{\vec{\pmb{\a}},p_1\}$ stands for $\{\a_1,\cdots,\a_k,p_1\}$ if $\vec{\pmb{\a}}=\{\a_1,\cdots,\a_k\}$. The symbol $\shuffle$ means summing over all
possible shuffles of two ordered sets $\vec{\pmb{\b}}_1$ and $\vec{\pmb{\b}}_2$, i.e., all permutations in the set $\vec{\pmb{\b}}_1\cup \vec{\pmb{\b}}_2$ while preserving the orderings
of $\vec{\pmb{\b}}_1$ and $\vec{\pmb{\b}}_2$. For instance, suppose $\vec{\pmb{\b}}_1=\{1,2\}$ and $\vec{\pmb{\b}}_2=\{3,4\}$, then
%
\bea
\vec{\pmb{\b}}_1\shuffle \vec{\pmb{\b}}_2=\{1,2,3,4\}+\{1,3,2,4\}+\{1,3,4,2\}+\{3,1,2,4\}+\{3,1,4,2\}+\{3,4,1,2\}\,.~~~~\label{shuffle}
\eea
%
The summation in \eref{expan-YMS} is over all ordered subsets $\vec{\pmb{\a}}=\{\a_1,\cdots,\a_k\}$, and the tensors $F_{\vec{\pmb{\a}}}$ are defined by
%
\bea
F^{\mu\nu}_{\vec{\pmb{\a}}}=\Big(f_{\a_k}\cdot f_{\a_{k-1}}\cdots f_{\a_2}\cdot f_{\a_1}\Big)^{\mu\nu}\,,
\eea
%
with
%
\bea
f^{\mu\nu}_{a}\equiv k^\mu_a\epsilon^\nu_a-\epsilon^\mu_a k^\nu_a\,,~~\label{defin-f}
\eea
%
where $k_a$ and $\epsilon_a$ are momentum and polarization carried by the external leg $a$, respectively. The combinatory momentum $Y_{\vec{\pmb{\a}}}$ is defined as $Y_{\vec{\pmb{\a}}}=Y_{\a_1}$ where $\a_1$ is the first element in the ordered set $\vec{\pmb{\a}}$, and $Y_{\a_1}$ is defined as the summation of momenta carried by external scalars at the l.h.s of $\a_1$ in the color ordering. For example, for ${\cal A}_{\rm YS}(1,2,\a_1,\cdots;\cdots|\sigma_{n+m})$, we have $Y_{\a_1}=k_1+k_2$.
In the expansion \eref{expan-YMS}, the coefficients are independent of the second color ordering $\sigma_{n+m}$, due to the well known double copy structure.  Using the expanded formula \eref{expan-YMS} iteratively, one can expand any single trace YMS amplitude to pure bi-adjoint scalar (BAS) amplitudes.

In the expansion \eref{expan-YMS}, the external gluon $p_1$ is special. Of course one can chose any $p_j$ with $j\in\{2,\cdots,m\}$ to instead of $p_1$, but a particular gluon is always required. Suppose $a$ is an element of the set $\pmb{\a}$, the coefficient $(\epsilon_{p_1}\cdot F_{\vec{\pmb{\a}}}\cdot Y_{\a_1})$ in \eref{expan-YMS} vanishing under the replacement $\epsilon_a\to k_a$, due to the definition of $f_a^{\mu\nu}$ in \eref{defin-f}. This character reflects the gauge invariance for the polarization $\epsilon_a$.
However, the gauge invariance for the polarization $\epsilon_1$ is obscure. If one use \eref{expan-YMS} iteratively to expand the YMS amplitude to BAS ones, in the resulting expansion, non of polarizations carries the manifest gauge invariance, since each external gluon will play the role of the special gluon once in the iterative procedure. After expanding to BAS amplitudes, the basis only includes propagators, thus the gauge invariance for each polarization is completely determined by coefficients. Thus, there must be some way to construct the coefficients which are manifestly gauge invariant. The above consideration leads to a natural question, how to achieve a new expansion which manifests the gauge invariance for each polarization of external gluons?

Such new expansion was first found by Clifford Cheung and James Mangan in \cite{Cheung:2021zvb}, via the so called covariant color-kinematic duality framework. Their result is equivalent to the following formula\footnote{The expansion in \cite{Cheung:2021zvb} is for expanding single trace Einstein-Yang-Mills (EYM) amplitudes to pure Yang-Mills (YM) ones, but the well known double copy structure \cite{Kawai:1985xq,Bern:2008qj,Chiodaroli:2014xia,Johansson:2015oia,Johansson:2019dnu}indicates that the expansion of YMS amplitudes in \eref{expan-YMS-gi} shares the same coefficients. The expanded formula in \cite{Cheung:2021zvb} is not the same as in \eref{expan-YMS-gi}, however it is straightforward to show the equivalence between them.}
%
\bea
& &{\cal A}_{\rm YS}(1,\cdots,n;p_1,\cdots,p_m|\sigma_{n+m})\nn
&=&\sum_{\vec{\pmb{\a}}}\,{k_b\cdot F_{\vec{\pmb{\a}}}\cdot Y_{\vec{\pmb{\a}}}\over k_b\cdot k_{p_1\cdots p_m}}\,{\cal A}_{\rm YS}(1,\{2,\cdots,n-1\}\shuffle\vec{\pmb{\a}},n;\{p_1,\cdots,p_m\}\setminus{\pmb{\a}}|\sigma_{n+m})\,,~~\label{expan-YMS-gi}
\eea
%
where $\pmb{\a}$ are un-empty subsets of $\{p_1,\cdots,p_m\}$. Here $k_b$ is a reference momentum, and the combinatory momentum $k_{p_1\cdots p_m}$ is defined as
$k_{p_1\cdots p_m}\equiv \sum_{i=1}^m k_{p_i}$. In \eref{expan-YMS-gi}, the gauge invariance for each polarization is manifested, with the cost of introducing the spurious pole $1/k_b\cdot k_{p_1\cdots p_m}$. Using \eref{expan-YMS-gi} iteratively to expand the YMS amplitude to pure BAS ones, the manifest gauge invariance for each polarization will be kept, although the manifest locality will be spoiled since various spurious poles will occur.

The method in \cite{Cheung:2021zvb} is based on the traditional Lagrangian and equations of motion. In the view of the experience of modern developments of scattering amplitudes, it is interesting to ask whether the expansion \eref{expan-YMS-gi} can be obtained only through on-shell imformation? This question is the main motivation of the current short note. In this note, we reproduce the expansion \eref{expan-YMS-gi} with the special choice of reference momentum $k_b=k_n$, from a totally different point of view. We first use the well known Bern-Carrasco-Johansson (BCJ) relation \cite{Bern:2008qj,Chiodaroli:2014xia,Johansson:2015oia,Johansson:2019dnu} to turn the expansion \eref{expan-YMS} to \eref{expan-YMS-gi}, for the special case that the single trace YMS amplitude includes only one external gluon. Then, we use the soft bootstrap method to generate the expansions of single trace YMS amplitudes with more external gluons. More explicitly, we consider the following soft theorem for the gluon at the sub-leading order \cite{Casali:2014xpa,Schwab:2014xua,Afkhami-Jeddi:2014fia},
%
\bea
{\cal A}^{(1)}_{\rm YS}(1,\cdots,n;p_1,\cdots,p_{m+1}|\sigma_{n+m+1})=S^{(1)}_g(p_{m+1})\,{\cal A}_{\rm YS}(1,\cdots,n;p_1,\cdots,p_m|\sigma_{n+m+1}\setminus p_{m+1})\,,~~\label{soft-theo}
\eea
%
where $S^{(1)}_g(p_{m+1})$ is the sub-leading soft operator. Here the momentum carried by the external gluon $p_{m+1}$
is re-scaled as $k_{p_{m+1}}\to\tau k_{p_{m+1}}$, and the amplitude is expanded in $\tau$. One can think the object ${\cal A}^{(1)}_{\rm YS}(1,\cdots,n;p_1,\cdots,p_{m+1}|\sigma_{n+m+1})$ as the sub-leading contribution in such Laurent expansion. If the expansion of
${\cal A}_{\rm YS}(1,\cdots,n;p_1,\cdots,p_m|\sigma_{n+m+1}\setminus p_{m+1})$ is already known, then the sub-leading soft theorem in \eref{soft-theo} indicates the expansion of ${\cal A}^{(1)}_{\rm YS}(1,\cdots,n;p_1,\cdots,p_{m+1}|\sigma_{n+m+1})$, therefore allows us to construct the expansions of YMS amplitudes with more external gluons recursively.

The remainder of this note is organized as follows. In section.\ref{sec-background}, we rapidly introduce the necessary background including tree BAS amplitudes, expansions of amplitudes to BAS ones, as well as soft theorems for external scalar and gluon. In section.\ref{sec-expan-YMS}, we derive the expansion \eref{expan-YMS-gi} via the soft bootstrap method. Then, we end with a brief summery in section.\ref{sec-conclusion}.

%%%%%%%%%%%%%%%%%%%%%%%%%%%%%%%%%%%%%%%%
\section{Background}
\label{sec-background}
%%%%%%%%%%%%%%%%%%%%%%%%%%%%%%%%%%%%%%%%

For readers' convenience, in this section we rapidly review the necessary background. In subsection.\ref{subsecBAS}, we review the tree level amplitudes of bi-adjoint scalar (BAS) theory. Some notations and technics which will be used in subsequent sections are also introduced. In subsection.\ref{subsecexpand}, we discuss the expansions of tree amplitudes to BAS amplitudes, including the choice of basis, as well as the double copy structure for coefficients. Finally, in subsection.\ref{subsec-soft}, we review the soft theorems for external scalars and gluons.

%%%%%%%%%%%%%%%%%%%%%%%%%%%%%%%%%%%%
\subsection{Tree level BAS amplitudes}
\label{subsecBAS}
%%%%%%%%%%%%%%%%%%%%%%%%%%%%%%%%%%%%%%%

The BAS theory describes the bi-adjoint scalar field $\phi_{a\bar{a}}$ with the Lagrangian
%
\bea
{\cal L}_{\rm BAS}={1\over2}\,\partial_\mu\phi^{a\bar{a}}\,\partial^{\mu}\phi^{a\bar{a}}-{\lambda\over3!}\,f^{abc}f^{\bar{a}\bar{b}\bar{c}}\,
\phi^{a\bar{a}}\phi^{b\bar{b}}\phi^{c\bar{c}}\,,
\eea
%
where the structure constant $f^{abc}$ and generator $T^a$ satisfy
%
\bea
[T^a,T^b]=if^{abc}T^c\,,
\eea
%
and the dual color algebra encoded by $f^{\bar{a}\bar{b}\bar{c}}$ and $T^{\bar{a}}$ is analogous.
The tree level amplitudes of this theory contain only propagators, and can be decomposed into double color ordered partial amplitudes via the standard technic.
Each double color ordered partial amplitude is simultaneously planar with respect to two color orderings, arise from expanding the full $n$-point amplitude to $\Tr(T^{a_{\sigma_1}}\cdots T^{a_{\sigma_n}})$ and $\Tr(T^{\bar{a}_{\bar{\sigma}_1}}\cdots T^{\bar{a}_{\bar{\sigma}_n}})$ respectively,
where $\sigma_i$ and $\bar{\sigma}_i$ denote permutations among all external scalars.

To calculate double color ordered partial amplitudes, it is convenient to employ the diagrammatical method proposed by Cachazo, He and Yuan in \cite{Cachazo:2013iea}.
For a BAS amplitude whose double color-orderings are given, this method provides the corresponding Feynman diagrams as well as the overall sign directly.
Let us consider the $5$-point example ${\cal A}_{\rm S}(1,2,3,4,5|1,4,2,3,5)$.
In Figure.\ref{5p}, the first diagram satisfies both two color orderings $(1,2,3,4,5)$ and $(1,4,2,3,5)$, while the second one satisfies the ordering
$(1,2,3,4,5)$ but not $(1,4,2,3,5)$. Thus, the first diagram is allowed by the double color orderings $(1,2,3,4,5|1,4,2,3,5)$, while the second one is not. It is easy to see other diagrams are also forbidden by the ordering $(1,4,2,3,5)$, thus the first diagram in Figure.\ref{5p} is the only diagram contributes to the amplitude
${\cal A}_{\rm S}(1,2,3,4,5|1,4,2,3,5)$.
\begin{figure}
  \centering
  % Requires \usepackage{graphicx}
  \includegraphics[width=6cm]{5p1.pdf}
   \includegraphics[width=6cm]{5p2.pdf}  \\
  \caption{Two $5$-point diagrams}\label{5p}
\end{figure}
The Feynman diagrams contribute to a given BAS amplitude can be obtained via systematic diagrammatical rules. For the above example, one can draw a disk diagram as follows.
%
\begin{itemize}
\item Draw points on the boundary of the disk according to the first ordering $(1,2,3,4,5)$.
\item Draw a loop of line segments which connecting the points according to the second ordering $(1,4,2,3,5)$.
\end{itemize}
%
The obtained disk diagram is shown in the first diagram in Figure.\ref{dis14235}. From the diagram, one can see that two orderings share the boundaries $\{1,5\}$ and $\{2,3\}$. These co-boundaries
indicate channels ${1/s_{15}}$ and ${1/s_{23}}$, therefore the first Feynman diagram in Figure.\ref{5p}. Then the BAS amplitude ${\cal A}_{\rm S}(1,2,3,4,5|1,4,2,3,5)$ can be computed as
%
\bea
{\cal A}_{\rm S}(1,2,3,4,5|1,4,2,3,5)={1\over s_{23}}{1\over s_{51}}\,,
\eea
%
up to an overall sign. The Mandelstam variable $s_{i\cdots j}$ is defined as
%
\bea
s_{i\cdots j}\equiv k_{i\cdots j}^2\,,~~~~k_{i\cdots j}\equiv\sum_{a=i}^j\,k_a\,,~~~~\label{mandelstam}
\eea
%
where $k_a$ is the momentum carried by the external leg $a$.

\begin{figure}
  \centering
  % Requires \usepackage{graphicx}
   \includegraphics[width=4cm]{14235.pdf}
   \includegraphics[width=4cm]{12435.pdf} \\
  \caption{Diagram for ${\cal A}_{\rm S}(1,2,3,4,5|1,4,2,3,5)$ and ${\cal A}_{\rm S}(1,2,3,4,5|1,2,4,3,5)$.}\label{dis14235}
\end{figure}



As another example, let us consider the BAS amplitude ${\cal A}_{\rm S}(1,2,3,4,5|1,2,4,3,5)$. The corresponding disk diagram is shown in the second configuration in
Figure.\ref{dis14235}, and one can see two orderings have co-boundaries $\{3,4\}$ and $\{5,1,2\}$. The co-boundary $\{3,4\}$ indicates the channel ${1/ s_{34}}$. The co-boundary $\{5,1,2\}$ indicates the channel ${1/s_{512}}$ which is equivalent to $1/s_{34}$, as well as sub-channels ${1/ s_{12}}$ and ${1/ s_{51}}$. Using the above decomposition, one can calculate ${\cal A}_{\rm S}(1,2,3,4,5|1,2,4,3,5)$ as
%
\bea
{\cal A}_{\rm S}(1,2,3,4,5|1,2,4,3,5)={1\over s_{34}}\Big({1\over s_{12}}+{1\over s_{51}}\Big)\,,
\eea
%
up to an overall sign.


The overall sign, determined by the color algebra, can be fixed by following rules.
%
\begin{itemize}
\item Each polygon with odd number of vertices contributes
a plus sign if its orientation is the same as that of the disk and a minus sign if opposite.
\item Each polygon with even number of vertices always contributes a minus sign.
\item Each intersection point contributes a minus sign.
\end{itemize}
%
We now apply these rules to previous examples. In the first diagram in Figure.\ref{dis14235}, the polygons are three triangles, namely $51A$, $A4B$ and $B23$, which contribute $+$, $-$, $+$ respectively, while two intersection points $A$ and $B$ contribute two $-$. In the second one in Figure.\ref{dis14235}, the polygons are $512A$ and $A43$, which contribute two $-$, while the intersection point $A$ contributes $-$. Then we arrive at the full results
%
\bea
{\cal A}_{\rm S}(1,2,3,4,5|1,4,2,3,5)&=&-{1\over s_{23}}{1\over s_{51}}\,,\nn
{\cal A}_{\rm S}(1,2,3,4,5|1,2,4,3,5)&=&-{1\over s_{34}}\Big({1\over s_{12}}+{1\over s_{51}}\Big)\,.
\eea


In the reminder of this paper, we adopt another convention for the overall sign. If the line segments form a convex polygon, and the orientation of the convex polygon is the same as that of the disk, then the overall sign is $+$. For instance, the disk diagram in Figure.\ref{newconvention} indicates the overall sign $+$ under the new convention, while the old convention gives $-$ according to the square formed by four line segments. Notice that the diagrammatical rules described previously still give the related sign between different disk diagrams. For example, two disk diagrams in Figure.\ref{dis14235} shows that the relative sign between ${\cal A}_{\rm S}(1,2,3,4,5|1,4,2,3,5)$ and
${\cal A}_{\rm S}(1,2,3,4,5|1,2,4,3,5)$ is $+$. The advantage of the new convention is that when removing a soft external scalar, the resulting sub-amplitude carries the same sign as the original one.

\begin{figure}
  \centering
  % Requires \usepackage{graphicx}
   \includegraphics[width=4cm]{newconvention.pdf} \\
  \caption{The overall sign $+$ under the new convention.}\label{newconvention}
\end{figure}

When considering the soft limit, the $2$-point channels play the central role. Since the partial BAS amplitude carries two color orderings, if the $2$-point channel contributes $1/s_{ab}$ to the amplitude, external legs $a$ and $b$ must be adjacent to each other in both two orderings.
Suppose the first color ordering is $(\cdots,a,b,\cdots)$, then $1/s_{ab}$ is allowed by this ordering. To denote if it is allowed by another one, we introduce the symbol $\delta_{ab}$ whose ordering of two subscripts $a$ and $b$ is determined by the first color ordering\footnote{The Kronecker symbol will not appear in this paper, thus we hope the notation $\delta_{ab}$ will not confuse the readers.}. The value of $\delta_{ab}$ is $\delta_{ab}=1$ if another color ordering is $(\cdots,a,b,\cdots)$, $\delta_{ab}=-1$ if another color ordering is $(\cdots,b,a,\cdots)$, due to the ani-symmetry of the structure constant, i.e., $f^{abc}=-f^{bac}$, and $\delta_{ab}=0$
otherwise. From the definition, it is straightforward to see $\delta_{ab}=-\delta_{ba}$. The symbol $\delta_{ab}$ will be used frequently latter.


%%%%%%%%%%%%%%%%%%%%%%%%%%%%%%%%%%%%%
\subsection{Expanding tree level amplitudes to BAS basis}
\label{subsecexpand}
%%%%%%%%%%%%%%%%%%%%%%%%%%%%%%%%%%%%%%%%%%

Tree level amplitudes which contain only massless particles can be expanded to double color ordered BAS amplitudes,
due to the observation that each Feynman diagram for pure propagators can be mapped to at least one disk diagram whose polygons are all triangles. An illustrative example is given in Figure.\ref{example}. Since each tree amplitude can be expanded by tree Feynman diagrams, and each Feynman diagram contributes propagators together with a numerator without any pole, one can conclude that each tree amplitude can be expanded to double color ordered BAS amplitudes, with coefficients which are polynomials depend on Lorentz invariants created by external kinematical variables.
To realize the expansion, one need to find the basis consists of BAS amplitudes.
Such basis can be determined by the well known Kleiss-Kuijf (KK) relation \cite{Kleiss:1988ne}
%
\bea
{\cal A}_{\rm S}(1,\vec{\pmb{\a}},n,\vec{\pmb{\b}}|\sigma_n)=(-)^{|\pmb{\a}|}\,{\cal A}_{\rm S}(1,\vec{\pmb{\a}}\shuffle\vec{\pmb{\b}}^T,n|\sigma_n)\,.~~~\label{KK}
\eea
%
Here $\vec{\pmb{\a}}$ and $\vec{\pmb{\b}}$ are two ordered subsets of external scalars, and $\vec{\pmb{\b}}^T$ stands for the ordered set generated from $\vec{\pmb{\b}}$ by reversing the original ordering. The $n$-point BAS amplitude ${\cal A}_{\rm S}(1,\vec{\pmb{\a}},n,\vec{\pmb{\b}}|\sigma_n)$ at the l.h.s of \eref{KK} carries two color orderings, one is $(1,\vec{\pmb{\a}},n,\vec{\pmb{\b}})$, another one is denoted by $\sigma_n$. The symbol $\shuffle$ means summing over all
possible shuffles of two ordered subsets, as explained in section.\ref{sec-intro}.
The analogous KK relation holds for another color ordering $\sigma_n$.
The KK relation implies that different double color ordered BAS amplitudes are not independent, thus the basis can be chosen as BAS amplitudes
${\cal A}_{\rm S}(1,\sigma_1,n|1,\sigma_2,n)$, with $1$ and $n$ are fixed at two ends in each color ordering. We call such basis the KK BAS basis. Based on the discussion above, the KK BAS basis can provide any structure of propagators, thus any amplitude can be expanded to this basis\footnote{The well known Bern-Carrasco-Johansson (BCJ) relation \cite{Bern:2008qj,Chiodaroli:2014xia,Johansson:2015oia,Johansson:2019dnu} links BAS amplitudes in the KK basis together, and the independent BAS amplitudes can be obtained by fixing three legs at three particular positions in the color orderings. However, in the BCJ relation, coefficients of BAS amplitudes depend on Mandelstam variables, this character leads to poles in coefficients when expanding to BCJ basis. On the other hand, when expanding to KK basis, one can find the expanded formula in which the coefficients contain no pole. In this paper, the coefficients of BAS amplitudes in KK basis contain spurious poles caused by manifesting the gauge invariance, but all physical poles are included in BAS amplitudes.}.

\begin{figure}
  \centering
  % Requires \usepackage{graphicx}
   \includegraphics[width=4cm]{example1.pdf}
   \includegraphics[width=4cm]{example2.pdf} \\
  \caption{Map between Feynman diagram and disk diagram.}\label{example}
\end{figure}

In this paper, we will consider the expansion of single trace YMS amplitudes. The single trace YMS amplitude ${\cal A}_{\rm YS}(1,\sigma_1,n;p_1,\cdots,p_m|1,\sigma_2,n)$ can be expanded to KK BAS basis as
%
\bea
& &{\cal A}_{\rm YS}(1,\sigma_{n-2},n;p_1,\cdots,p_m|1,\sigma_{n+m-2},n)\nn
&=&\sum_{\sigma'_{n+m-2}}\,{\cal C}(\sigma_{n-2},\sigma'_{n+m-2},\epsilon_i,k_i)\,{\cal A}_{\rm S}(1,\sigma'_{n+m-2},n|1,\sigma_{n+m-2},n)\,,~~~\label{exp-YS-KK}
\eea
%
where $\sigma'_{n+m-2}$ are permutations among external legs in $\{2,\cdots,n-1\}\cup\{p_1,\cdots,p_m\}$. The double copy structure \cite{Kawai:1985xq,Bern:2008qj,Chiodaroli:2014xia,Johansson:2015oia,Johansson:2019dnu} indicates that the coefficient ${\cal C}(\sigma_{n-2},\sigma'_{n+m-2},\epsilon_i,k_i)$ depends on polarization vectors $\epsilon_i$ of external gluons,
momenta $k_i$ of all external particles, permutations $\sigma'_{n+m-2}$ and $\sigma_{n-2}$, but is independent of the permutation $\sigma_{n+m-2}$\footnote{Originally, the double copy means the GR amplitude can be factorized as ${\cal A}_{\rm G}={\cal A}_{\rm Y}\times {\cal S}\times{\cal A}_{\rm Y}$, where the kernel ${\cal S}$ is obtained by inverting propagators. Our assumption that the coefficients depend on only one color ordering is equivalent to the original version, see in \cite{Zhou:2022orv}.}. Thus, suppose we replace $(1,\sigma_{n+m-2},n)$ by the more general ordering $\sigma_{n+m}$ among all external legs, without fixing $1$ and $n$ at any position, the expansion in \eref{exp-YS-KK} still holds.
The construction of coefficients ${\cal C}(\sigma_{n-2},\sigma'_{n+m-2},\epsilon_i,k_i)$ is the main goal of this note.


%%%%%%%%%%%%%%%%%%%%%%%%%%%%%%%%%%%
\subsection{Soft theorems for external scalars and gluons}
\label{subsec-soft}
%%%%%%%%%%%%%%%%%%%%%%%%%%%%%%%%%%%%%

In this subsection, we introduce the soft theorems for external scalars and gluons, which are crucial for subsequent discussions in the next subsection.

Consider the double color ordered BAS amplitude ${\cal A}_{\rm S}(1,\cdots,n|\sigma_n)$. We re-scale $k_i$
as $k_i\to\tau k_i$, and expand the amplitude in $\tau$. The leading order contribution manifestly aries from $2$-point channels $1/s_{1(i+1)}$
and $1/ s_{(i-1)i}$ which provide the $1/\tau$ order contributions, namely,
%
\bea
{\cal A}^{(0)}_{\rm S}(1,\cdots,n|\sigma_n)&=&{1\over \tau}\Big({\delta_{i(i+1)}\over s_{i(i+1)}}+{\delta_{(i-1)i}\over s_{(i-1)i}}\Big)\,
{\cal A}_{\rm S}(1,\cdots,i-1,\not{i},i+1,\cdots,n|\sigma_n\setminus i)\nn
&=&S^{(0)}_s(i)\,{\cal A}_{\rm S}(1,\cdots,i-1,\not{i},i+1,\cdots,n|\sigma_n\setminus i)\,,~~~\label{soft-theo-s}
\eea
%
where $\not{i}$ stands for removing the leg $i$, $\sigma_n\setminus1$ means the color ordering generated from $\sigma_n$ by eliminating $i$. The leading soft operator $S^{(0)}_s(i)$ for the scalar $i$ is extracted as
%
\bea
S^{(0)}_s(i)={1\over \tau}\,\Big({\delta_{i(i+1)}\over s_{i(i+1)}}+{\delta_{(i-1)i}\over s_{(i-1)i}}\Big)\,,~~~~\label{soft-fac-s-0}
\eea
%
which acts on external scalars which are adjacent to $i$ in two color orderings.
By imposing the universality of soft factor, the above result can be generated to the YMS amplitude as
%
\bea
{\cal A}^{(0)}_{\rm YS}(1,\cdots,n;p_1,\cdots,p_m|\sigma_{n+m})
&=&S^{(0)}_s(i)\,{\cal A}_{\rm YS}(1,\cdots,i-1,\not{i},i+1,\cdots,n;p_1,\cdots,p_m|\sigma_{n+m}\setminus i)\,,~~~\label{soft-theo-s}
\eea
%
where the soft factor $S^{(0)}_s(i)$ is the same as in \eref{soft-fac-s-0}. In other words, the soft operator $S^{(0)}_s(i)$ does not act on any external gluon.

The soft theorems for external gluons at leading and sub-leading orders can be obtained via various approaches \cite{Casali:2014xpa,Schwab:2014xua,Afkhami-Jeddi:2014fia,Zhou:2022orv}. Notice that one of these methods is to use the expanded formula of YMS amplitude in \eref{expan-YMS}, as can be seen in \cite{Zhou:2022orv}. Such soft theorems are found to be
%
\bea
{\cal A}^{(0)}_{\rm YS}(1,\cdots,n;p_1,\cdots,p_m|\sigma_{n+m})
&=&S^{(0)}_g(p_i)\,{\cal A}_{\rm YS}(1,\cdots,n;p_1,\cdots,\not{p_i},\cdots,p_m|\sigma_{n+m}\setminus p_i)\,,~~~\label{soft-theo-g0}
\eea
%
and
%
\bea
{\cal A}^{(0)}_{\rm YS}(1,\cdots,n;p_1,\cdots,p_m|\sigma_{n+m})
&=&S^{(1)}_g(p_i)\,{\cal A}_{\rm YS}(1,\cdots,n;p_1,\cdots,\not{p_i},\cdots,p_m|\sigma_{n+m}\setminus p_i)\,,~~~\label{soft-theo-g1}
\eea
%
where the external momentum $k_{p_i}$ is re-scaled as $k_{p_i}\to\tau k_{p_i}$. The soft factors at two orders are given by
%
\bea
S^{(0)}_g(p_i)={1\over\tau}\,\sum_{a\in\sigma}\,{\delta_{ap_i}\,(\epsilon_{p_i}\cdot k_a)\over s_{ap_i}}\,,~~\label{soft-fac-g-0}
\eea
%
and
%
\bea
S^{(1)}_g(p_i)=\sum_{a\in\sigma}\,{\delta_{ap_1}\,\big(\epsilon_{p_i}\cdot J_a\cdot k_{p_i}\big)\over s_{ap_1}}\,,~~\label{soft-fac-g-1}
\eea
%
respectively. In \eref{soft-fac-g-0} and \eref{soft-fac-g-1}, $a$ should be summed over all external legs, i.e., these soft operators for external gluon act on both external scalars and gluons.

The sub-leading soft operator \eref{soft-fac-g-1} for external gluon will play the central role in the next subsection. Here we mention some useful results when acting this operator to Lorentz vectors. The angular momentum operator $J_a^{\mu\nu}$ acts on Lorentz vector $k^\rho_a$ with the orbital part of the generator and on $\epsilon^\rho_a$ with the spin part of the generator in the vector representation as follows,
%
\bea
J_a^{\mu\nu}\,k_a^\rho=k_a^{[\mu}\,{\partial k_a^\rho\over\partial k_{a,\nu]}}\,,~~~~
J_a^{\mu\nu}\,\epsilon_a^\rho=\big(\eta^{\nu\rho}\,\delta^\mu_\sigma-\eta^{\mu\rho}\,\delta^\nu_\sigma\big)\,\epsilon^\sigma_a\,.
\eea
%
These actions can be summarized as
%
\bea
S^{(1)}_g(p)&=&-\sum_{V_a}\,{\delta_{ap}\over s_{ap}}\,V_a\cdot f_p\cdot{\partial\over\partial V_a}\,,~~\label{act-s-V}
\eea
%
due to the observation that the amplitude is linear in each polarization vector. In \eref{act-s-V}, the summation over $V_a$
is among all Lorentz vectors including both momenta and polarizations. One can understood \eref{act-s-V} as a differential operator which satisfies the Leibnitz's rule. Using \eref{act-s-V}, we immediately get
%
\bea
\big(S^{(1)}_g(p)\,k_a\big)\cdot V=-{\delta_{ap}\over s_{ap}}\,(k_a\cdot f_p\cdot V)\,,~~~~\big(S^{(1)}_g(p)\,\epsilon_a\big)\cdot V=-{\delta_{ap}\over s_{ap}}\,(\epsilon_a\cdot f_p\cdot V)\,,~~~~\label{iden-1}
\eea
%
where $V$ is an arbitrary Lorentz vector, and
%
\bea
V_1\cdot\big(S^{(1)}_g(p)\,f_a\big)\cdot V_2={\delta_{ap}\over s_{ap}}\,V_1\cdot(f_p\cdot f_a-f_a\cdot f_p)\cdot V_2\,,~~~~\label{iden-2}
\eea
%
for two arbitrary Lorentz vectors $V_1$ and $V_2$, where the anti-symmetric tensor $f_a$ is defined as $f_a^{\mu\nu}\equiv k_a^\mu\epsilon_a^\nu-\epsilon_a^\mu k_a^\nu$.

%%%%%%%%%%%%%%%%%%%%%%%%%%%%%%%%%%%%
\section{Expansion of single trace YMS amplitude}
\label{sec-expan-YMS}
%%%%%%%%%%%%%%%%%%%%%%%%%%%%%%%%%%%%%

In this section, we derive the expansion of single trace YMS amplitude to BAS KK basis, in the formula which is manifestly gauge invariant for each polarization. We first use the well known Bern-Carrasco-Johansson (BCJ) relation to rewrite the expansion \eref{expan-YMS} for the special BAS amplitudes with only one external gluon into a manifestly gauge invariant form, in subsection.\ref{subsec-1gluon}. Then, we use a method based on the sub-leading soft theorem for external gluons, to give a recursive pattern which leads to the general expansion of single trace YMS amplitudes with arbitrary number of external gluons,
which is manifestly gauge invariant for each polarization.

%%%%%%%%%%%%%%%%%%%%%%%%%%%%%%%%%%%%%%%
\subsection{YMS amplitude with one external gluon}
\label{subsec-1gluon}
%%%%%%%%%%%%%%%%%%%%%%%%%%%%%%%%%%%%%%%

We start by considering the single trace YMS amplitude ${\cal A}_{\rm YS}(1,\cdots,n;p|\sigma_{n+1})$ which contains external scalars $i\in\{1,\cdots,n\}$
and an external gluon $p$. As can be seen in \cite{Zhou:2022orv}, the soft theorem for scalars, together with the universality of soft factor, completely fixes the expanded formula of
YMS amplitude ${\cal A}_{\rm YS}(1,\cdots,n;p|\sigma_{n+1})$ to be,
%
\bea
{\cal A}_{\rm YS}(1,\cdots,n;p|\sigma_{n+1})=(\epsilon_p\cdot Y_p)\,{\cal A}_{\rm S}(1,\{2,\cdots,n-1\}\shuffle p,n|\sigma_{n+1})\,,~~~~\label{expan-1g-old}
\eea
%
where the combinatory momentum $Y_p$ was defined in section.\ref{sec-intro} after \eref{defin-f}.
In the above expansion, it is direct to observe the gauge invariance of the polarization $\epsilon_p$, since the replacement $\epsilon_p^\mu\to k_p^\mu$ yields the well known BCJ relation
%
\bea
0=(k_p\cdot Y_p)\,{\cal A}_{\rm S}(1,\{2,\cdots,n-1\}\shuffle p,n|\sigma_{n+1})\,.
\eea
%
However, if one use the formula \eref{expan-1g-old} and the soft theorem for external gluon to derive the expansions of YMS amplitudes with more external gluons, the resulting expansions are in the formula in \eref{expan-YMS} \cite{Zhou:2022orv}, and
the gauge invariance for the polarization of the special gluon is obscure. The proof of gauge invariance for such special polarization requires the application of BCJ relations in a complicated way, and the complexity increases rapidly as the number of external gluons increases. To avoid this disadvantage, we want to
rewrite \eref{expan-1g-old} so that the manifest gauge invariance is carried by coefficients. To realize the goal,
one can use the BCJ relation to modify \eref{expan-1g-old} as
%
\bea
{\cal A}_{\rm YS}(1,\cdots,n;p|\sigma_{n+1})
&=&(\epsilon_p\cdot Y_p)\,{\cal A}_{\rm S}(1,\{2,\cdots,n-1\}\shuffle p,n|\sigma_{n+1})\nn
& &-{k_n\cdot\epsilon_p\over k_n\cdot k_p}\,(k_p\cdot Y_p)\,{\cal A}_{\rm S}(1,\{2,\cdots,n-1\}\shuffle p,n|\sigma_{n+1})\nn
&=&{k_n\cdot f_p\cdot Y_p\over k_n\cdot k_p}\,{\cal A}_{\rm S}(1,\{2,\cdots,n-1\}\shuffle p,n|\sigma_{n+1})\,.~~~~\label{expan-1g-gi}
\eea
%
In the last line at the r.h.s of \eref{expan-1g-gi}, the coefficients vanish under the replacement $\epsilon_p^\mu\to k_p^\mu$,
due to the definition of the tensor $f^{\mu\nu}_p\equiv k^\mu_p\epsilon^\nu_p-\epsilon^\mu_p k^\nu_p$. This is the desired new expansion for the YMS with one external gluon, trivialized the gauge invariance of polarization $\epsilon^\mu_p$, with the cost that a spurious pole $k_n\cdot k_p$ is included in coefficients. As will be seen, by applying our method based on the soft theorem for external gluons, this new expansion leads to general expansion of YMS amplitudes with arbitrary number of external gluons, which manifests the gauge invariance for each polarization.

%%%%%%%%%%%%%%%%%%%%%%%%%%%%%%%%%%%%%%%
\subsection{YMS amplitude with two external gluons}
\label{subsec-2gluon}
%%%%%%%%%%%%%%%%%%%%%%%%%%%%%%%%%%%%%%%

Now we turn to the single trace YMS amplitude ${\cal A}_{\rm YS}(1,\cdots,n;p_1,p_2|\sigma_{n+2})$ with two external gluons labeled as $p_1$ and $p_2$. Let us
consider the soft behavior of gluon $p_2$, i.e., we re-scale $k_{p_2}$ as $k_{p_2}\to\tau k_{p_2}$, and expand ${\cal A}_{\rm YS}(1,\cdots,n;p_1,p_2|\sigma_{n+2})$ in $\tau$. The soft theorem \eref{soft-theo-g1} indicates that the contribution at sub-leading order should be
%
\bea
{\cal A}^{(1)}_{\rm YS}(1,\cdots,n;p_1,p_2|\sigma_{n+2})&=&S^{(1)}_g(p_2)\,{\cal A}_{\rm YS}(1,\cdots,n;p_1|\sigma_{n+2}\setminus p_2)\nn
&=&S^{(1)}_g(p_2)\,\Big[{k_n\cdot f_{p_1}\cdot Y_{p_1}\over k_n\cdot k_{p_1}}\,{\cal A}_{\rm S}(1,\{2,\cdots,n-1\}\shuffle p_1,n|\sigma_{n+2}\setminus p_2)\Big]\nn
&=&B_1+B_2+B_3\,,~~~\label{2g-Lb}
\eea
%
where
%
\bea
B_1={k_n\cdot f_{p_1}\cdot Y_{p_1}\over k_n\cdot k_{p_1}}\,{\cal A}^{(1)}_{\rm YS}(1,\{2,\cdots,n-1\}\shuffle p_1,n;p_2|\sigma_{n+2})\,,~~\label{B1}
\eea
%
%
\bea
B_2&=&\tau\,{k_n\cdot f_{p_2}\cdot f_{p_1}\cdot Y_{p_1}\over k_n\cdot k_{p_1}}\,{\cal A}^{(0)}_{\rm S}(1,\{2,\cdots,n-1\}\shuffle \{p_1,p_2\},n|\sigma_{n+2})\nn
& &+\tau\,{k_n\cdot f_{p_1}\cdot f_{p_2}\cdot Y_{p_2}\over k_n\cdot k_{p_1}}\,{\cal A}^{(0)}_{\rm S}(1,\{2,\cdots,n-1\}\shuffle \{p_2,p_1\},n|\sigma_{n+2})\,,~~\label{B2}
\eea
%
and
%
\bea
B_3=-\tau\,{(k_n\cdot f_{p_1}\cdot Y_{p_1})\,(k_n\cdot f_{p_2}\cdot k_{p_1})\over (k_n\cdot k_{p_1})^2}\,{\cal A}^{(0)}_{\rm S}(1,\{2,\cdots,n-1\}\shuffle \{p_1,p_2\},n|\sigma_{n+2})\,.~~~\label{B3}
\eea
%
The second equality of \eref{2g-Lb} is obtained by substituting the expansion \eref{expan-1g-gi}. Three parts $B_1$, $B_2$ and $B_3$ arise from
acting $S^{(1)}_g(p_2)$ on ${\cal A}_{\rm S}(1,\{2,\cdots,n-1\}\shuffle p_1,n|\sigma_{n+2}\setminus p_2)$, $k_n\cdot f_{p_1}\cdot Y_{p_1}$ and $k_n\cdot k_{p_1}$ respectively, due to Leibnitz's rule. The expression for the first part $B_1$ in \eref{B1} is obtained via the soft theorem
%
\bea
S^{(1)}_g(p_2)\,{\cal A}_{\rm S}(1,\{2,\cdots,n-1\}\shuffle p_1,n|\sigma_{n+2}\setminus p_2)={\cal A}^{(1)}_{\rm YS}(1,\{2,\cdots,n-1\}\shuffle p_1,n;p_2|\sigma_{n+2})\,.
\eea
%
The second part $B_2$ can be calculated as follows.
Using relations \eref{iden-1} and \eref{iden-2}, we have
%
\bea
S^{(1)}_g(p_2)\,(k_n\cdot f_{p_1}\cdot Y_{p_1})=H_1+H_2\,,
\eea
%
where
%
\bea
H_1&=&-{\delta_{np_2}\,(k_n\cdot f_{p_2}\cdot f_{p_1}\cdot Y_{p_1})\over s_{np_2}}
+{\delta_{p_1p_2}\,(k_n\cdot f_{p_2}\cdot f_{p_1}\cdot X_{p_1})\over s_{p_1p_2}}\nn
H_2&=&-{\delta_{p_1p_2}\,(k_n\cdot f_{p_1}\cdot f_{p_2}\cdot Y_{p_1})\over s_{p_1p_2}}
+\sum_{i=1}^j\,{\delta_{ip_2}\,(k_n\cdot f_{p_1}\cdot f_{p_2}\cdot k_i)\over s_{ip_2}}\,,~~~~\label{H1H2}
\eea
%
with $Y_{p_1}=\sum_{i=1}^j\,k_i$. We use $\delta_{ab}=-\delta_{ba}$ to reorganize $H_1$ as
%
\bea
H_1=\Big[\Big({\delta_{p_1p_2}\over s_{p_1p_2}}+{\delta_{p_2(j+1)}\over s_{p_2(j+1)}}\Big)+\sum_{k=j+1}^{n-1}\,\Big({\delta_{kp_2}\over s_{kp_2}}+{\delta_{p_2(k+1)}\over s_{p_2(k+1)}}\Big)\Big]\,(k_n\cdot f_{p_2}\cdot f_{p_1}\cdot Y_{p_1})\,.~~~~\label{h1}
\eea
%
The reason for expressing $H_1$ in the above manner is that we want to interpret $H_1\,{\cal A}_{\rm S}(1,\{2,\cdots,n-1\}\shuffle p_1,n|\sigma_{n+2}\setminus p_2)$ as coefficients times leading or sub-leading terms of amplitudes, as can be seen in \eref{B2}.
Combining \eref{h1} and the leading soft factor \eref{soft-fac-s-0} for the scalar, we find
%
\bea
& &H_1\,{\cal A}_{\rm S}(1,\{2,\cdots,n-1\}\shuffle p_1,n|\sigma_{n+2}\setminus p_2)\nn
&=&\sum_{j=1}^{n-1}\,H_1\,{\cal A}_{\rm S}(1,\cdots,j,p_1,j+1,\cdots,n|\sigma_{n+2}\setminus p_2)\nn
&=&\tau\,(k_n\cdot f_{p_2}\cdot f_{p_1}\cdot Y_{p_1})\,{\cal A}^{(0)}_{\rm S}(1,\{2,\cdots,n-1\}\shuffle\{p_1,p_2\},n|\sigma_{n+2})\,.~~\label{H1}
\eea
%
We emphasize that ${\cal A}^{(0)}_{\rm S}(1,\{2,\cdots,n-1\}\shuffle\{p_1,p_2\},n|\sigma_{n+2})$ is the leading order contribution of ${\cal A}_{\rm S}(1,\{2,\cdots,n-1\}\shuffle\{p_1,p_2\},n|\sigma_{n+2})$ which is at the $\tau^{-1}$ order, since $f_{p_2}$ in the coefficients is accompanied with $\tau$, while ${\cal A}^{(1)}_{\rm YS}(1,\cdots,n;p_1,p_2|\sigma_{n+2})$ under consideration is at the $\tau^0$ order. Based on the reason similar as that for rewriting $H_1$, we reorganize $H_2$ as
%
\bea
H_2&=&\sum_{i=1}^j\,\Big({\delta_{ip_2}\over s_{ip_2}}-{\delta_{p_1p_2}\over s_{p_1p_2}}\Big)\,(k_n\cdot f_{p_1}\cdot f_{p_2}\cdot k_i)\nn
&=&\sum_{i=1}^j\,\Big[\sum_{l=i}^{j-1}\,\Big({\delta_{lp_2}\over s_{lp_2}}+{\delta_{p_2(l+1)}\over s_{p_2(l+1)}}\Big)+{\delta_{jp_2}\over s_{jp_2}}+{\delta_{p_2p_1}\over s_{p_2p_1}}\Big]\,(k_n\cdot f_{p_1}\cdot f_{p_2}\cdot k_i)\nn
&=&\sum_{l=1}^{j}\,\Big({\delta_{lp_2}\over s_{lp_2}}+{\delta_{p_2(l+1)}\over s_{p_2(l+1)}}\Big)\,\Big[k_n\cdot f_{p_1}\cdot f_{p_2}\cdot\Big(\sum_{i=1}^l\,k_i\Big)\Big]\,,
\eea
%
where $\delta_{q(j+1)}$ and $s_{q(j+1)}$ in the last line should be understood as $j+1=p_1$. Thus,
%
\bea
& &H_2\,{\cal A}_{\rm S}(1,\{2,\cdots,n-1\}\shuffle p_1,n|\sigma_{n+2}\setminus p_2)\nn
&=&\sum_{j=1}^{n-1}\,\sum_{l=1}^{j}\,\Big({\delta_{lp_2}\over s_{lp_2}}+{\delta_{p_2(l+1)}\over s_{p_2(l+1)}}\Big)\,\Big[k_n\cdot f_{p_1}\cdot f_{p_2}\cdot\Big(\sum_{i=1}^l\,k_i\Big)\Big]\,{\cal A}_{\rm S}(1,\cdots,j,p_1,j+1,\cdots,n|\sigma_{n+2}\setminus p_2)\nn
&=&\sum_{j=1}^{n-1}\,\sum_{l=1}^{j}\,\tau\,\Big[k_n\cdot f_{p_1}\cdot f_{p_2}\cdot\Big(\sum_{i=1}^l\,k_i\Big)\Big]\,{\cal A}^{(0)}_{\rm S}(1,\cdots,l,p_2,l+1,\cdots,j,p_1,j+1,\cdots,n|\sigma_{n+2})\nn
&=&\tau\,(k_n\cdot f_{p_1}\cdot f_{p_2}\cdot Y_{p_2})\,{\cal A}^{(0)}_{\rm S}(1,\{2,\cdots,n-1\}\shuffle\{p_2,p_1\},n|\sigma_{n+2})\,.~~\label{H2}
\eea
%
Putting \eref{H1} and \eref{H2} together yields the expression of $B_2$ in \eref{B2}. The third part $B_3$ can be calculated by using the relation \eref{iden-1},
%
\bea
S^{(1)}_g(p_2)\,{1\over k_n\cdot k_{p_1}}&=&-{k_n\cdot f_{p_2}\cdot k_{p_1}\over(k_n\cdot k_{p_1})^2}\,\Big({\delta_{p_1p_2}\over s_{p_1p_2}}-{\delta_{np_2}\over s_{np_2}}\Big)\nn
&=&-{k_n\cdot f_{p_2}\cdot k_{p_1}\over(k_n\cdot k_{p_1})^2}\,\Big[{\delta_{p_1p_2}\over s_{p_1p_2}}+{\delta_{p_2a}\over s_{p_2a}}+\sum_{k=a}^{n-1}\,\Big({\delta_{kp_2}\over s_{kp_2}}+{\delta_{p_2(k+1)}\over s_{p_2(k+1)}}\Big)\Big]\,,
\eea
%
which indicates
%
\bea
& &\sum_{i=1}^{n-1}\,\Big[S^{(1)}_g(p_2)\,{1\over k_n\cdot k_{p_1}}\Big]\,{\cal A}_{\rm S}(1,\cdots,i,p_1,i+1,\cdots,n|\sigma_{n+2}\setminus p_2)\nn
&=&-\tau\,{k_n\cdot f_{p_2}\cdot k_{p_1}\over(k_n\cdot k_{p_1})^2}\,{\cal A}^{(0)}_{\rm S}(1,\{2,\cdots,n-1\}\shuffle\{p_1,p_2\},n|\sigma_{n+2})\,,
\eea
%
therefore we get the expression of $B_3$ in \eref{B3}.

We will use $B_1$, $B_2$ and $B_3$ to restore the complete expansion of ${\cal A}_{\rm YS}(1,\cdots,n;p_1,p_2|\sigma_{n+2})$. Mathematically, it is impossible to construct the full ${\cal A}_{\rm YS}(1,\cdots,n;p_1,p_2|\sigma_{n+2})$ from the sub-leading contribution ${\cal A}^{(1)}_{\rm YS}(1,\cdots,n;p_1,p_2|\sigma_{n+2})$, thus further physical constrains are needed. The most important one is the symmetry under the permutation $p_1\leftrightarrow p_2$. The physical amplitude ${\cal A}_{\rm YS}(1,\cdots,n;p_1,p_2|\sigma_{n+2})$ should be invariant under such re-labeling. In the expansion \eref{expan-YMS}, since a special gluon $p_1$ is required, such symmetry is not manifest. Instead, one can use the BCJ relations to prove that the resulting expansions obtained by choosing different special gluons  are equivalent to each other. However, the purpose of the current work is to seek the new expansion which manifests the gauge invariance for all polarizations, which means the democracy among all external gluons, thus it is natural to expect the symmetry under $p_1\leftrightarrow p_2$.

With the additional condition of symmetry $p_1\leftrightarrow p_2$, let us try to find the expansion of the amplitude ${\cal A}_{\rm YS}(1,\cdots,n;p_1,p_2|\sigma_{n+2})$. It is straightforward to recognize that $B_2$ arises as the leading term of
%
\bea
P_1&=&{k_n\cdot f_{p_2}\cdot f_{p_1}\cdot Y_{p_1}\over k_n\cdot k_{p_1p_2}}\,{\cal A}_{\rm S}(1,\{2,\cdots,n-1\}\shuffle \{p_1,p_2\},n|\sigma_{n+2})\nn
& &+{k_n\cdot f_{p_1}\cdot f_{p_2}\cdot Y_{p_2}\over k_n\cdot k_{p_1p_2}}\,{\cal A}_{\rm S}(1,\{2,\cdots,n-1\}\shuffle \{p_2,p_1\},n|\sigma_{n+2})\,,~~~~\label{2g-P1}
\eea
%
which is at the $\tau^0$ order. The reason for turning $k_n\cdot k_{p_1}$ in the denominators to $k_n\cdot k_{p_1p_2}$ is to keep the symmetry $p_1\leftrightarrow p_2$. There is another way to keep such symmetry, which is adding terms with $k_n\cdot k_{p_2}$ in the denominator. However, ${\cal A}^{(1)}_{\rm YS}(1,\cdots,n;p_1,p_2|\sigma_{n+2})$ shall receive contributions from such terms, for example the coefficients carry $1/(\tau k_n\cdot k_{p_2})$ under the re-scaling $k_{p_2}\to\tau k_{p_2}$, and the amplitudes contribute $\tau^1$, which combine to terms at the $\tau^0$ order. Since such contributions do not exist in ${\cal A}^{(1)}_{\rm YS}(1,\cdots,n;p_1,p_2|\sigma_{n+2})$ in \eref{2g-Lb}, the $1/k_n\cdot k_{p_2}$ terms are forbidden. Then the symmetry $p_1\leftrightarrow p_2$ also excludes the $1/k_n\cdot k_{p_1}$ terms, and the remaining choice is only $1/k_n\cdot k_{p_1p_2}$.

One can also observe that $B_1$ comes from the sub-leading term of
%
\bea
P_2={k_n\cdot f_{p_1}\cdot Y_{p_1}\over k_n\cdot k_{p_1p_2}}\,{\cal A}_{\rm YS}(1,\{2,\cdots,n-1\}\shuffle p_1,n;p_2|\sigma_{n+2})\,,~~\label{2g-P2}
\eea
%
which is also at the $\tau^0$ order. Again, the denominator is set to be $k_n\cdot k_{p_1p_2}$ due to the symmetry. However, the sub-leading contribution of $P_2$ receives additional contribution from expanding the denominator in $\tau$. Further more, the symmetry among two gluons requires the existence of
%
\bea
P_3={k_n\cdot f_{p_2}\cdot Y_{p_2}\over k_n\cdot k_{p_1p_2}}\,{\cal A}_{\rm YS}(1,\{2,\cdots,n-1\}\shuffle p_2,n;p_1|\sigma_{n+2})\,,~~\label{2g-P3}
\eea
%
which serves as the counter part of $P_2$ under the permutation $p_1\leftrightarrow p_2$.
These missing terms will be figured out by considering $B_3$. The expression of $B_3$ in \eref{B3} also breaks the symmetry among $p_1$ and $p_2$.
Based on such symmetry, the appearance of ${\cal A}^{(0)}_{\rm S}(1,\{2,\cdots,n-1\}\shuffle \{p_1,p_2\},n|\sigma_{n+2})$ in \eref{B3} implies the existence of ${\cal A}^{(0)}_{\rm S}(1,\{2,\cdots,n-1\}\shuffle \{p_2,p_1\},n|\sigma_{n+2})$. To realize this, we observe that
%
\bea
& &-(k_n\cdot f_{p_2}\cdot k_{p_1})\,{\cal A}^{(0)}_{\rm S}(1,\{2,\cdots,n-1\}\shuffle\{p_1,p_2\},n|\sigma_{n+2})\nn
&=&\Big(k_n\cdot f_{p_2}\cdot Y_{p_2}-k_n\cdot f_{p_2}\cdot X_{p_2} \Big)\,{\cal A}^{(0)}_{\rm S}(1,\{2,\cdots,n-1\}\shuffle p_1\shuffle p_2,n|\sigma_{n+2})\,,~~\label{reform-B3}
\eea
%
where the combinatory momentum $X_a$ is defined as the summation of momenta carried by external legs at the l.h.s of $a$ in the color ordering, without distinguishing scalars and gluons.
Then, $B_3$ can be identified as
%
\bea
B_3=B_4+B_5\,,
\eea
%
where
%
\bea
B_4=\tau\,{k_n\cdot f_{p_2}\cdot Y_{p_2}\over k_n\cdot k_{p_1}}\,{\cal A}^{(0)}_{\rm YS}(1,\{2,\cdots,n-1\}\shuffle p_2,n;p_1|\sigma_{n+2})\,,~~\label{B4}
\eea
%
and
%
\bea
B_5=-\tau\,{(k_n\cdot k_{p_2})\,(k_n\cdot f_{p_1}\cdot Y_{p_1})\over (k_n\cdot k_{p_1})^2}\,{\cal A}^{(0)}_{\rm YS}(1,\{2,\cdots,n-1\}\shuffle p_1,n;p_2|\sigma_{n+2})\,.~~\label{B5}
\eea
%
Here we have used expansions
%
\bea
{\cal A}_{\rm YS}(1,\{2,\cdots,n-1\}\shuffle p_2,n;p_1|\sigma_{n+2})&=&{k_n\cdot f_{p_1}\cdot Y_{p_1}\over k_n\cdot k_{p_1}}\,{\cal A}_{\rm S}(1,\{2,\cdots,n-1\}\shuffle p_1\shuffle p_2,n|\sigma_{n+2})\,,\nn
{\cal A}_{\rm YS}(1,\{2,\cdots,n-1\}\shuffle p_1,n;p_2|\sigma_{n+2})&=&{k_n\cdot f_{p_2}\cdot Y_{p_2}\over k_n\cdot k_{p_2}}\,{\cal A}_{\rm S}(1,\{2,\cdots,n-1\}\shuffle p_1\shuffle p_2,n|\sigma_{n+2})\,.~~\label{expansions}
\eea
%
Notice that $Y_{p_2}$ in \eref{expansions} equals to $Y_{p_2}$ in \eref{reform-B3} since in \eref{expansions} the leg $p_1$ is a scalar. The combinatory momentum $Y_{p_1}$ equals to $Y_{p_1}$ in \eref{B3} and \eref{expansions} at the leading order, since $k_{p_2}$ is accompanied with $\tau$ thus is dropped into the sub-leading term.

Now one can observe that $B_4$ is the leading term of $P_3$, and $B_5$ arises from the term at $\tau^1$ order in the expansion of the denominator of $P_2$. Thus, we conclude that
%
\bea
{\cal A}_{\rm YS}(1,\cdots,n;p_1,p_2|\sigma_{n+2})&=&P_1+P_2+P_3\nn
&=&{k_n\cdot f_{p_1}\cdot Y_{p_1}\over k_n\cdot k_{p_1p_2}}\,{\cal A}_{\rm YS}(1,\{2,\cdots,n-1\}\shuffle p_1,n;p_2|\sigma_{n+2})\nn
& &+{k_n\cdot f_{p_2}\cdot Y_{p_2}\over k_n\cdot k_{p_1p_2}}\,{\cal A}_{\rm YS}(1,\{2,\cdots,n-1\}\shuffle p_2,n;p_1|\sigma_{n+2})\nn
& &+{k_n\cdot f_{p_2}\cdot f_{p_1}\cdot Y_{p_1}\over k_n\cdot k_{p_1p_2}}\,{\cal A}_{\rm S}(1,\{2,\cdots,n-1\}\shuffle \{p_1,p_2\},n|\sigma_{n+2})\nn
& &+{k_n\cdot f_{p_1}\cdot f_{p_2}\cdot Y_{p_2}\over k_n\cdot k_{p_1p_2}}\,{\cal A}_{\rm S}(1,\{2,\cdots,n-1\}\shuffle \{p_2,p_1\},n|\sigma_{n+2})\,.~~\label{expan-2g-gi}
\eea
%
One can verify that taking $k_{p_1}\to\tau k_{p_1}$ and expanding \eref{expan-2g-gi} in $\tau$ give the correct behavior at sub-leading order, namely,
%
\bea
{\cal A}^{(1)}_{\rm YS}(1,\cdots,n;p_1,p_2|\sigma_{n+2})=S^{(1)}_g(p_1)\,{\cal A}_{\rm YS}(1,\cdots,n;p_2|\sigma_{n+2}\setminus p_1)\,.
\eea
%
Notice that in the formula \eref{expan-2g-gi}, both the symmetry among two gluons and the gauge invariance for each polarization are manifested.

%%%%%%%%%%%%%%%%%%%%%%%%
\subsection{YMS amplitude with three external gluons}
\label{subsec-3gluon}
%%%%%%%%%%%%%%%%%%%%%%%%%%%%%%

The next example is the single trace YMS amplitude ${\cal A}_{\rm YS}(1,\cdots,n;p_1,p_2,p_3|\sigma_{n+3})$ which includes three external gluons $p_1$, $p_2$ and $p_3$. The treatment is parallel to that in the previous subsection for the $2$-gluon case, and such similarity exhibits the recursive pattern. Similar as in the previous case, we consider
$k_{p_3}\to\tau k_{p_3}$, and expand ${\cal A}_{\rm YS}(1,\cdots,n;p_1,p_2,p_3|\sigma_{n+3})$ in $\tau$. The soft theorem requires the sub-leading term to be
%
\bea
& &{\cal A}^{(1)}_{\rm YS}(1,\cdots,n;p_1,p_2,p_3|\sigma_{n+3})\nn&=&S^{(1)}_g(p_3)\,{\cal A}_{\rm YS}(1,\cdots,n;p_1,p_2|\sigma_{n+3}\setminus p_3)\nn
&=&S^{(1)}_g(p_3)\,\Big[{k_n\cdot f_{p_1}\cdot Y_{p_1}\over k_n\cdot k_{p_1p_2}}\,{\cal A}_{\rm YS}(1,\{2,\cdots,n-1\}\shuffle p_1,n;p_2|\sigma_{n+3}\setminus p_3)+p_1\leftrightarrow p_2\nn
& &+{k_n\cdot f_{p_2}\cdot f_{p_1}\cdot Y_{p_1}\over k_n\cdot k_{p_1p_2}}\,{\cal A}_{\rm S}(1,\{2,\cdots,n-1\}\shuffle \{p_1,p_2\},n|\sigma_{n+3}\setminus p_3)+p_1\leftrightarrow p_2\Big]\,,
\eea
%
where the expanded formula \eref{expan-2g-gi} for YMS amplitude with two external gluons is used. Applying the Leibnitz's rule, we get
%
\bea
{\cal A}^{(1)}_{\rm YS}(1,\cdots,n;p_1,p_2,p_3|\sigma_{n+3})=B_1+B_2+B_3+B_4+B_5+B_6\,,
\eea
%
where
%
\bea
B_1&=&{k_n\cdot f_{p_1}\cdot Y_{p_1}\over k_n\cdot k_{p_1p_2}}\,\Big[S^{(1)}_g(p_3)\,{\cal A}_{\rm YS}(1,\{2,\cdots,n-1\}\shuffle p_1,n;p_2|\sigma_{n+3}\setminus p_3)\Big]+p_1\leftrightarrow p_2\nn
&=&{k_n\cdot f_{p_1}\cdot Y_{p_1}\over k_n\cdot k_{p_1p_2}}\,{\cal A}^{(1)}_{\rm YS}(1,\{2,\cdots,n-1\}\shuffle p_1,n;p_2,p_3|\sigma_{n+3})+p_1\leftrightarrow p_2\,,~~\label{3g-B1}
\eea
%
%
\bea
B_2&=&{k_n\cdot f_{p_2}\cdot f_{p_1}\cdot Y_{p_1}\over k_n\cdot k_{p_1p_2}}\,\Big[S^{(1)}_g(p_3)\,{\cal A}_{\rm S}(1,\{2,\cdots,n-1\}\shuffle \{p_1,p_2\},n|\sigma_{n+3}\setminus p_3)\Big]+p_1\leftrightarrow p_2\nn
&=&{k_n\cdot f_{p_2}\cdot f_{p_1}\cdot Y_{p_1}\over k_n\cdot k_{p_1p_2}}\,{\cal A}^{(1)}_{\rm YS}(1,\{2,\cdots,n-1\}\shuffle \{p_1,p_2\},n;p_3|\sigma_{n+3})+p_1\leftrightarrow p_2\,,~~\label{3g-B2}
\eea
%
%
\bea
B_3&=&{\Big[S^{(1)}_g(p_3)\,k_n\cdot f_{p_1}\cdot Y_{p_1}\Big]\over k_n\cdot k_{p_1p_2}}\,{\cal A}_{\rm YS}(1,\{2,\cdots,n-1\}\shuffle p_1,n;p_2|\sigma_{n+3}\setminus p_3)+p_1\leftrightarrow p_2\nn
&=&\tau\,{k_n\cdot f_{p_3}\cdot f_{p_1}\cdot Y_{p_1}\over k_n\cdot k_{p_1p_2}}\,{\cal A}^{(0)}_{\rm YS}(1,\{2,\cdots,n-1\}\shuffle \{p_1,p_3\},n;p_2|\sigma_{n+3})+p_1\leftrightarrow p_2\nn
& &+\tau\,{k_n\cdot f_{p_1}\cdot f_{p_3}\cdot Y_{p_3}\over k_n\cdot k_{p_1p_2}}\,{\cal A}^{(0)}_{\rm YS}(1,\{2,\cdots,n-1\}\shuffle \{p_3,p_1\},n;p_2|\sigma_{n+3})+p_1\leftrightarrow p_2\,,~~\label{3g-B3}
\eea
%
%
\bea
B_4&=&{\Big[S^{(1)}_g(p_3)\,k_n\cdot f_{p_2}\cdot f_{p_1}\cdot Y_{p_1}\Big]\over k_n\cdot k_{p_1p_2}}\,{\cal A}_{\rm S}(1,\{2,\cdots,n-1\}\shuffle \{p_1,p_2\},n|\sigma_{n+3}\setminus p_3)+p_1\leftrightarrow p_2\nn
&=&\tau\,{k_n\cdot f_{p_3}\cdot f_{p_2}\cdot f_{p_1}\cdot Y_{p_1}\over k_n\cdot k_{p_1p_2}}\,{\cal A}^{(0)}_{\rm S}(1,\{2,\cdots,n-1\}\shuffle \{p_1,p_2,p_3\},n|\sigma_{n+3})+p_1\leftrightarrow p_2\nn
& &+\tau\,{k_n\cdot f_{p_2}\cdot f_{p_3}\cdot f_{p_1}\cdot Y_{p_1}\over k_n\cdot k_{p_1p_2}}\,{\cal A}^{(0)}_{\rm S}(1,\{2,\cdots,n-1\}\shuffle \{p_1,p_3,p_2\},n|\sigma_{n+3})+p_1\leftrightarrow p_2\nn
& &+\tau\,{k_n\cdot f_{p_2}\cdot f_{p_1}\cdot f_{p_3}\cdot Y_{p_3}\over k_n\cdot k_{p_1p_2}}\,{\cal A}^{(0)}_{\rm S}(1,\{2,\cdots,n-1\}\shuffle \{p_3,p_1,p_2\},n|\sigma_{n+3})+p_1\leftrightarrow p_2\,,~~\label{3g-B4}
\eea
%
%
\bea
B_5&=&\Big[S^{(1)}_g(p_3)\,{1\over k_n\cdot k_{p_1p_2}}\Big]\,(k_n\cdot f_{p_1}\cdot Y_{p_1})\,{\cal A}_{\rm YS}(1,\{2,\cdots,n-1\}\shuffle p_1,n;p_2|\sigma_{n+3}\setminus p_3)+p_1\leftrightarrow p_2\nn
&=&-\tau\,{(k_n\cdot f_{p_1}\cdot Y_{p_1})\,(k_n\cdot f_{p_3}\cdot K_{p_3})\over \big(k_n\cdot k_{p_1p_2}\big)^2}\,\W{\cal A}^{(0)}_{\rm YS}(1,\{2,\cdots,n-1\}\shuffle p_1\shuffle p_3,n;p_2|\sigma_{n+3})+p_1\leftrightarrow p_2\,,~~\label{3g-B5}
\eea
%
%
\bea
B_6&=&\Big[S^{(1)}_g(p_3)\,{1\over k_n\cdot k_{p_1p_2}}\Big]\,(k_n\cdot f_{p_2}\cdot f_{p_1}\cdot Y_{p_1})\,{\cal A}_{\rm S}(1,\{2,\cdots,n-1\}\shuffle \{p_1,p_2\},n|\sigma_{n+3}\setminus p_3)+p_1\leftrightarrow p_2\nn
&=&-\tau\,{(k_n\cdot f_{p_2}\cdot f_{p_1}\cdot Y_{p_1})\,(k_n\cdot f_{p_3}\cdot K_{p_3})\over \big(k_n\cdot k_{p_1p_2} \big)^2}\,{\cal A}^{(0)}_{\rm S}(1,\{2,\cdots,n-1\}\shuffle \{p_1,p_2\}\shuffle p_3,n|\sigma_{n+3}\setminus p_3)\nn
& &+p_1\leftrightarrow p_2\,.~~\label{3g-B6}
\eea
%
The first and second parts $B_1$ and $B_2$ in \eref{3g-B1} and \eref{3g-B2} are analogous to $B_1$ in \eref{B1}, come from the soft theorem. The computation of $B_3$ and $B_4$ in \eref{3g-B3} and \eref{3g-B4}
are parallel to that for obtaining $B_2$ in \eref{B2}. The calculation of $B_5$ and $B_6$ in \eref{3g-B5} and \eref{3g-B6} are parallel to the process for obtaining $B_3$ in \eref{B3}. The combinatory momentum $K_{p_3}$ is defined as the summation of momenta in $\{k_{p_1},k_{p_2}\}$
whose corresponding legs are at the l.h.s of $p_3$ in the color ordering. The notation $\W{\cal A}^{(0)}_{\rm YS}(1,\{2,\cdots,n-1\}\shuffle p_1\shuffle p_3,n;p_2|\sigma_{n+3})$ in the last line of \eref{3g-B5} means expanding ${\cal A}^{(0)}_{\rm YS}(1,\{2,\cdots,n-1\}\shuffle p_1\shuffle p_3,n;p_2|\sigma_{n+3})$ to pure BAS basis, and $K_{p_3}$ are defined for the corresponding color orderings in these BAS amplitudes. If one use other expressions for ${\cal A}^{(0)}_{\rm YS}(1,\{2,\cdots,n-1\}\shuffle p_1\shuffle p_3,n;p_2|\sigma_{n+3})$, \eref{3g-B5} does not hold.

Using $K_{p_3}=Y_{p_3}-X_{p_3}$, where for \eref{3g-B5} $X_{p_3}$ is defined for BAS amplitudes after expanding ${\cal A}^{(0)}_{\rm YS}(1,\{2,\cdots,n-1\}\shuffle p_1\shuffle p_3,n;p_2|\sigma_{n+3})$, one can separate $B_5$ and $B_6$ as
%
\bea
B_5=B_7+B_8\,,~~~~B_6=B_9+B_{10}\,,
\eea
%
where
%
\bea
B_7=\tau\,{(k_n\cdot f_{p_1}\cdot Y_{p_1})\,(k_n\cdot f_{p_3}\cdot Y_{p_3})\over \big(k_n\cdot k_{p_1p_2}\big)^2}\,{\cal A}^{(0)}_{\rm YS}(1,\{2,\cdots,n-1\}\shuffle p_1\shuffle p_3,n;p_2|\sigma_{n+3})+p_1\leftrightarrow p_2\,,~~\label{3g-B7}
\eea
%
%
\bea
B_8=-\tau\,{(k_n\cdot k_{p_3})\,(k_n\cdot f_{p_1}\cdot Y_{p_1})\over \big(k_n\cdot k_{p_1p_2}\big)^2}\,{\cal A}^{(0)}_{\rm YS}(1,\{2,\cdots,n-1\}\shuffle p_1,n;p_2,p_3|\sigma_{n+3})+p_1\leftrightarrow p_2\,,~~\label{3g-B8}
\eea
%
%
\bea
B_9&=&\tau\,{(k_n\cdot f_{p_2}\cdot f_{p_1}\cdot Y_{p_1})\,(k_n\cdot f_{p_3}\cdot Y_{p_3})\over \big(k_n\cdot k_{p_1p_2}\big)^2}\,{\cal A}^{(0)}_{\rm S}(1,\{2,\cdots,n-1\}\shuffle \{p_1,p_2\}\shuffle p_3,n|\sigma_{n+3})\nn
& &+p_1\leftrightarrow p_2\,~~\label{3g-B9}
\eea
%
%
\bea
B_{10}=-\tau\,{(k_n\cdot k_{p_3})\,(k_n\cdot f_{p_2}\cdot f_{p_1}\cdot Y_{p_1})\over \big(k_n\cdot k_{p_1p_2}\big)^2}\,{\cal A}^{(0)}_{\rm YS}(1,\{2,\cdots,n-1\}\shuffle \{p_1,p_2\},n;p_3|\sigma_{n+3})+p_1\leftrightarrow p_2\,,~~\label{3g-B10}
\eea
%
One can obtain $B_8$ in \eref{3g-B8} as follows. Using the expansion of ${\cal A}_{\rm YS}(1,\{2,\cdots,n-1\}\shuffle p_1,n;p_2,p_3|\sigma_{n+3})$, we have
%
\bea
& &{\cal A}^{(0)}_{\rm YS}(1,\{2,\cdots,n-1\}\shuffle p_1,n;p_2,p_3|\sigma_{n+3})\nn
&=&{k_n\cdot f_{p_2}\cdot Y_{p_2}\over k_n\cdot k_{p_2}}\,{\cal A}^{(0)}_{\rm YS}(1,\{2,\cdots,n-1\}\shuffle p_1\shuffle p_2,n;p_3|\sigma_{n+3})\,,~~\label{B81}
\eea
%
since all other terms are at the $\tau^1$ order. Then we use the expansion of ${\cal A}_{\rm YS}(1,\{2,\cdots,n-1\}\shuffle p_1\shuffle p_2,n;p_3|\sigma_{n+3})$ to get
%
\bea
& &{\cal A}^{(0)}_{\rm YS}(1,\{2,\cdots,n-1\}\shuffle p_1\shuffle p_2,n;p_3|\sigma_{n+3})\nn
&=&{k_n\cdot f_{p_3}\cdot X_{p_3}\over k_n\cdot k_{p_3}}\,{\cal A}^{(0)}_{\rm S}(1,\{2,\cdots,n-1\}\shuffle p_1\shuffle p_2\shuffle p_3,n|\sigma_{n+3})\,.~~\label{B82}
\eea
%
Notice that for pure BAS amplitudes $Y_{p_3}$ is equivalent to $X_{p_3}$. Expanding ${\cal A}^{(0)}_{\rm YS}(1,\{2,\cdots,n-1\}\shuffle p_1\shuffle p_2,n;p_3|\sigma_{n+3})$ will not alter $Y_{p_2}$ in \eref{B81} at the leading order, since $k_{p_3}$
is accompanied with $\tau$ thus does not contribute. Substituting \eref{B82} into \eref{B81}, we find
%
\bea
& &{\cal A}^{(0)}_{\rm YS}(1,\{2,\cdots,n-1\}\shuffle p_1,n;p_2,p_3|\sigma_{n+3})\nn
&=&{k_n\cdot f_{p_3}\cdot X_{p_3}\over k_n\cdot k_{p_3}}\,\W{\cal A}^{(0)}_{\rm YS}(1,\{2,\cdots,n-1\}\shuffle p_1\shuffle p_3,n;p_2|\sigma_{n+3})\,,~~\label{B83}
\eea
%
where the expansion of ${\cal A}_{\rm YS}(1,\{2,\cdots,n-1\}\shuffle p_1\shuffle p_3,n;p_2|\sigma_{n+3})$ has been used. The relation \eref{B83} yields the expression of $B_8$ in \eref{3g-B8}. The expression of $B_9$ in \eref{3g-B10} can be calculated in the analogous way.

Now one can recognize that $B_1$ in \eref{3g-B1} together with $B_8$ in \eref{3g-B8} gives the sub-leading term of
%
\bea
P_1={k_n\cdot f_{p_1}\cdot Y_{p_1}\over k_n\cdot k_{p_1p_2p_3}}\,{\cal A}_{\rm YS}(1,\{2,\cdots,n-1\}\shuffle p_1,n;p_2,p_3|\sigma_{n+3})+p_1\leftrightarrow p_2\,,~~\label{3g-P1}
\eea
%
while $B_2$ in \eref{3g-B2} and $B_{10}$ in \eref{3g-B10} gives the sub-leading term of
%
\bea
P_2={k_n\cdot f_{p_2}\cdot f_{p_1}\cdot Y_{p_1}\over k_n\cdot k_{p_1p_2p_3}}\,{\cal A}_{\rm YS}(1,\{2,\cdots,n-1\}\shuffle \{p_1,p_2\},n;p_3|\sigma_{n+3})+p_1\leftrightarrow p_2\,.~~\label{3g-P2}
\eea
%
Similar as in the $2$-gluon case, $k_n\cdot k_{p_1p_2}$ in the denominator is turned to $k_{n}\cdot k_{p_1p_2p_3}$, to keep the manifest symmetry among external gluons $p_1$, $p_2$, $p_3$. $B_3$ in \eref{3g-B3} can be interpreted as the leading term of
%
\bea
P_3&=&{k_n\cdot f_{p_3}\cdot f_{p_1}\cdot Y_{p_1}\over k_n\cdot k_{p_1p_2p_3}}\,{\cal A}_{\rm YS}(1,\{2,\cdots,n-1\}\shuffle \{p_1,p_3\},n;p_2|\sigma_{n+3})+p_1\leftrightarrow p_2\nn
& &+{k_n\cdot f_{p_1}\cdot f_{p_3}\cdot Y_{p_3}\over k_n\cdot k_{p_1p_2p_3}}\,{\cal A}_{\rm YS}(1,\{2,\cdots,n-1\}\shuffle \{p_3,p_1\},n;p_2|\sigma_{n+3})+p_1\leftrightarrow p_2\,,~~\label{3g-B3}
\eea
%
$B_4$ in \eref{3g-B4} can be interpreted as the leading term of
%
\bea
P_4&=&{k_n\cdot f_{p_3}\cdot f_{p_2}\cdot f_{p_1}\cdot Y_{p_1}\over k_n\cdot k_{p_1p_2p_3}}\,{\cal A}_{\rm S}(1,\{2,\cdots,n-1\}\shuffle \{p_1,p_2,p_3\},n|\sigma_{n+3})+p_1\leftrightarrow p_2\nn
& &+{k_n\cdot f_{p_2}\cdot f_{p_3}\cdot f_{p_1}\cdot Y_{p_1}\over k_n\cdot k_{p_1p_2p_3}}\,{\cal A}_{\rm S}(1,\{2,\cdots,n-1\}\shuffle \{p_1,p_3,p_2\},n|\sigma_{n+3})+p_1\leftrightarrow p_2\nn
& &+{k_n\cdot f_{p_2}\cdot f_{p_1}\cdot f_{p_3}\cdot Y_{p_3}\over k_n\cdot k_{p_1p_2p_3}}\,{\cal A}_{\rm S}(1,\{2,\cdots,n-1\}\shuffle \{p_3,p_1,p_2\},n|\sigma_{n+3})+p_1\leftrightarrow p_2\,.~~\label{3g-P4}
\eea
%
Finally, combining $B_7$ in \eref{3g-B7} and $B_9$ in \eref{3g-B9} together gives the leading terms of
%
\bea
P_5={k_n\cdot f_{p_3}\cdot Y_{p_3}\over k_n\cdot k_{p_1p_2p_3}}\,{\cal A}_{\rm YS}(1,\{2,\cdots,n-1\}\shuffle p_3,n;p_1,p_2|\sigma_{n+3})\,.~~\label{3g-P5}
\eea
%
Thus, we arrive at
%
\bea
& &{\cal A}_{\rm YS}(1,\cdots,n;p_1,p_2,p_3|\sigma_{n+3})\nn
&=&P_1+P_2+P_3+P_4+P_5\nn
&=&\sum_{\a_1\in\{p_1,p_2,p_2\}}\,{k_n\cdot f_{\a_1}\cdot Y_{\a_1}\over k_n\cdot k_{p_1p_2p_3}}\,{\cal A}_{\rm YS}(1,\{2,\cdots,n-1\}\shuffle \a_1,n;\{p_1,p_2,p_3\}\setminus\a_1|\sigma_{n+3})\nn
& &+\sum_{\b_1,\b_2\in\{p_1,p_2,p_3\}}\,{k_n\cdot f_{\b_2}\cdot f_{\b_1}\cdot Y_{\b_1}\over k_n\cdot k_{p_1p_2p_3}}\,{\cal A}_{\rm YS}(1,\{2,\cdots,n-1\}\shuffle \{\b_1,\b_2\},n;\{p_1,p_2,p_3\}\setminus\{\b_1,\b_2\}|\sigma_{n+3})\nn
& &+\sum_{\gamma_i\in\{p_1,p_2,p_3\}}\,{k_n\cdot f_{\gamma_3}\cdot f_{\gamma_2}\cdot f_{\gamma_1}\cdot Y_{\b_1}\over k_n\cdot k_{p_1p_2p_3}}\,{\cal A}_{\rm S}(1,\{2,\cdots,n-1\}\shuffle \{\gamma_1,\gamma_2,\gamma_3\},n|\sigma_{n+3})\,,~~\label{expan-3g-gi}
\eea
%
where $i\in\{1,2,3\}$ in the final line.

%%%%%%%%%%%%%%%%%%%%%%%%%%%%%
\subsection{General case}
\label{subsec-general}
%%%%%%%%%%%%%%%%%%%%%%%%%%%%%%%

Comparing the manipulations in subsections.\ref{subsec-2gluon} and \ref{subsec-3gluon}, we see the processes are extremely parallel to each other. This similarity implies the recursive pattern which yields the general expansion for single trace YMS amplitudes with arbitrary number of external gluons. Such general expansion is the goal of this subsection.

In general, the single trace YMS amplitude ${\cal A}_{\rm YS}(1,\cdots,n;p_1,\cdots,p_m|\sigma_{n+m})$ with $m$ external gluons can be expanded as
in \eref{expan-YMS-gi}. When $\pmb{\a}=\{p_1,\cdots,p_m\}$, the YMS amplitudes at the second line of \eref{expan-YMS-gi} are reduced to pure BAS ones. As mentioned in section.\ref{sec-intro}, in this note we restrict ourselves to the special choice of the reference momentum $k_b$ in \eref{expan-YMS-gi}, which is $k_b=k_n$.

Obviously, expansions in \eref{expan-1g-gi}, \eref{expan-2g-gi}, \eref{expan-3g-gi} satisfy the general formula in \eref{expan-YMS-gi} with $k_b=k_n$. To get this general formula, let us prove that suppose \eref{expan-YMS-gi} is satisfied for a particular $m$, then it is also satisfied for $m+1$. Consider $k_{p_{m+1}}\to\tau k_{p_{m+1}}$, and expand ${\cal A}_{\rm YS}(1,\cdots,n;p_1,\cdots,p_m|\sigma_{n+m+1})$ in $\tau$ to get
%
\bea
{\cal A}^{(1)}_{\rm YS}(1,\cdots,n;p_1,\cdots,p_{m+1}|\sigma_{n+m+1})
&=&S^{(1)}_g(p_{m+1})\,{\cal A}_{\rm YS}(1,\cdots,n;p_1,\cdots,p_m|\sigma_{n+m+1}\setminus p_{m+1})\nn
&=&B_1+B_2+B_3\,,~~\label{3B}
\eea
%
where
%
\bea
B_1&=&\sum_{\vec{\pmb{\a}}}\,{k_n\cdot F_{\vec{\pmb{\a}}}\cdot Y_{\vec{\pmb{\a}}}\over k_n\cdot k_{p_1\cdots p_m}}\nn
& &\Big[S^{(1)}_g(p_{m+1}){\cal A}_{\rm YS}(1,\{2,\cdots,n-1\}\shuffle\vec{\pmb{\a}},n;\{p_1,\cdots,p_m\}\setminus\pmb{\a}|\sigma_{n+m+1}\setminus p_{m+1})\Big]\nn
&=&\sum_{\vec{\pmb{\a}}}\,{k_n\cdot F_{\vec{\pmb{\a}}}\cdot Y_{\vec{\pmb{\a}}}\over k_n\cdot k_{p_1\cdots p_m}}\nn
& &{\cal A}^{(1)}_{\rm YS}(1,\{2,\cdots,n-1\}\shuffle\vec{\pmb{\a}},n;\{\{p_1,\cdots,p_m\}\setminus\pmb{\a}\}\cup p_{m+1}|\sigma_{n+m+1})\,,~~\label{mg-B1}
\eea
%
%
\bea
B_2
&=&\sum_{\vec{\pmb{\a}}}\,{\Big[S^{(1)}_g(p_{m+1})\,k_n\cdot F_{\vec{\pmb{\a}}}\cdot Y_{\vec{\pmb{\a}}}\Big]\over k_n\cdot k_{p_1\cdots p_m}}\nn
& &{\cal A}_{\rm YS}(1,\{2,\cdots,n-1\}\shuffle\vec{\pmb{\a}},n;\{p_1,\cdots,p_m\}\setminus\pmb{\a}|\sigma_{n+m+1}\setminus p_{m+1})\nn
&=&\sum_{\vec{\pmb{\a}}'}\,\tau\,{k_n\cdot F_{\vec{\pmb{\a}}'}\cdot Y_{\vec{\pmb{\a}}'}\over k_n\cdot k_{p_1\cdots p_m}}\nn
& &{\cal A}^{(0)}_{\rm YS}(1,\{2,\cdots,n-1\}\shuffle\vec{\pmb{\a}}',n;\{p_1,\cdots,p_m,p_{m+1}\}\setminus\pmb{\a'}|\sigma_{n+m+1})\,,~~\label{mg-B2}
\eea
%
with $\vec{\pmb{\a}}'=\vec{\pmb{\a}}\shuffle p_{m+1}$, and
%
\bea
B_3
&=&\sum_{\vec{\pmb{\a}}}\,\Big[S^{(1)}_g(p_{m+1})\,{1\over k_n\cdot k_{p_1\cdots p_m}}\Big]\,(k_n\cdot F_{\vec{\pmb{\a}}}\cdot Y_{\vec{\pmb{\a}}})\nn
& &{\cal A}_{\rm YS}(1,\{2,\cdots,n-1\}\shuffle\vec{\pmb{\a}},n;\{p_1,\cdots,p_m\}\setminus\pmb{\a}|\sigma_{n+m+1}\setminus p_{m+1})\nn
&=&-\sum_{\vec{\pmb{\a}}}\,\tau\,{(k_n\cdot F_{\vec{\pmb{\a}}}\cdot Y_{\vec{\pmb{\a}}})\,(k_n\cdot f_{p_{m+1}}\cdot K_{p_{m+1}})\over \big(k_n\cdot k_{p_1\cdots p_m}\big)^2}\nn
& &\W{\cal A}^{(0)}_{\rm YS}(1,\{2,\cdots,n-1\}\shuffle\vec{\pmb{\a}}\shuffle p_{m+1},n;\{p_1,\cdots,p_m\}\setminus\pmb{\a}|\sigma_{n+m+1})\,.~~\label{mg-B3}
\eea
%
The calculation are paralleled to those in subsections \ref{subsec-2gluon} and \ref{subsec-3gluon}. The notation $\W{\cal A}^{(0)}_{\rm YS}(1,\{2,\cdots,n-1\}\shuffle\vec{\pmb{\a}}\shuffle p_{m+1},n;\{p_1,\cdots,p_m\}\setminus\pmb{\a}|\sigma_{n+m+1})$ is again understood as expanding ${\cal A}^{(0)}_{\rm YS}(1,\{2,\cdots,n-1\}\shuffle\vec{\pmb{\a}}\shuffle p_{m+1},n;\{p_1,\cdots,p_m\}\setminus\pmb{\a}|\sigma_{n+m+1})$ to BAS basis. To obtain $B_2$ in the second line of \eref{mg-B2}, we used the property
%
\bea
\vec{\pmb{a}}\shuffle\vec{\pmb{b}}\shuffle\vec{\pmb{c}}=\vec{\pmb{a}}\shuffle\{\vec{\pmb{b}}\shuffle\vec{\pmb{c}}\}\,,
\eea
%
where $\vec{\pmb{a}}$, $\vec{\pmb{b}}$ and $\vec{\pmb{c}}$ are three ordered sets. In \eref{mg-B3}, $K_{p_{m+1}}$ stands for the summation of momenta carried by external gluons $p_i$ at the l.h.s of $p_{m+1}$ in the color ordering. We emphasize that the external leg $p_{m+1}$ is not included in the set $\pmb{\a}$ appear
in $B_1$, $B_2$ and $B_3$.

Using $K_{p_{m+1}}=Y_{p_{m+1}}-X_{p_{m+1}}$, one can split $B_3$ as
%
\bea
B_3=B_4+B_5\,,
\eea
%
where
%
\bea
B_4&=&\sum_{\vec{\pmb{\a}}}\,\tau\,{(k_n\cdot F_{\vec{\pmb{\a}}}\cdot Y_{\vec{\pmb{\a}}})\,(k_n\cdot f_{p_{m+1}}\cdot Y_{p_{m+1}})\over \big(k_n\cdot k_{p_1\cdots p_m}\big)^2}\nn
& &{\cal A}^{(0)}_{\rm YS}(1,\{2,\cdots,n-1\}\shuffle\vec{\pmb{\a}}\shuffle p_{m+1},n;\{p_1,\cdots,p_m\}\setminus\pmb{\a}|\sigma_{n+m+1})\nn
&=&\tau\,{k_n\cdot f_{p_{m+1}}\cdot Y_{p_{m+1}}\over k_n\cdot k_{p_1\cdots p_m}}\,{\cal A}^{(0)}_{\rm YS}(1,\{2,\cdots,n-1\}\shuffle p_{m+1},n;p_1,\cdots,p_m|\sigma_{n+m+1})\,,~~\label{mg-B4}
\eea
%
%
\bea
B_5&=&-\sum_{\vec{\pmb{\a}}}\,\tau\,{(k_n\cdot F_{\vec{\pmb{\a}}}\cdot Y_{\vec{\pmb{\a}}})\,(k_n\cdot f_{p_{m+1}}\cdot X_{p_{m+1}})\over \big(k_n\cdot k_{p_1\cdots p_m}\big)^2}\nn
& &\W{\cal A}^{(0)}_{\rm YS}(1,\{2,\cdots,n-1\}\shuffle\vec{\pmb{\a}}\shuffle p_{m+1},n;\{p_1,\cdots,p_m\}\setminus\pmb{\a}|\sigma_{n+m+1})\nn
&=&-\sum_{\vec{\pmb{\a}}}\,\tau\,{(k_n\cdot k_{p_{m+1}})\,(k_n\cdot F_{\vec{\pmb{\a}}}\cdot Y_{\vec{\pmb{\a}}})\over \big(k_n\cdot k_{p_1\cdots p_m}\big)^2}\nn
& &{\cal A}^{(0)}_{\rm YS}(1,\{2,\cdots,n-1\}\shuffle\vec{\pmb{\a}},n;\{\{p_1,\cdots,p_m\}\setminus\pmb{\a}\}\cup p_{m+1}|\sigma_{n+m+1})\,.~~\label{mg-B5}
\eea
%
The computations of $B_4$ and $B_5$ are analogous to those for obtaining $B_7$, $B_8$, $B_9$ and $B_{10}$ in subsection. \ref{subsec-3gluon}. It is worth to give a brief discussion for how to generalize the manipulation from \eref{B81} to \eref{B83} to the current general case. We first observe that
%
\bea
& &{\cal A}^{(0)}_{\rm YS}(1,\{2,\cdots,n-1\}\shuffle\vec{\pmb{\a}},n;\{\{p_1,\cdots,p_m\}\setminus\pmb{\a}\}\cup p_{m+1}|\sigma_{n+m+1})\nn
&=&\sum_{\vec{\pmb{\b}}}\,{k_n\cdot F_{\vec{\pmb{\b}}}\cdot Y_{\vec{\pmb{\b}}}\over k_n\cdot k_{\{p_i\}\setminus\pmb{\a}}}\,
{\cal A}^{(0)}_{\rm YS}(1,\{2,\cdots,n-1\}\shuffle\vec{\pmb{\a}}\shuffle\vec{\pmb{\b}},n;\{\{p_1,\cdots,p_m\}\setminus\{\pmb{\a}\cup\pmb{\b}\}\}\cup p_{m+1}|\sigma_{n+m+1})\,.
~~\label{toB5-1}
\eea
%
Here $\pmb{\b}$ is a subset of $\{p_1,\cdots,p_m\}\setminus\pmb{\a}$ which does not include $p_{m+1}$, since otherwise $k_{p_{m+1}}$ will occur in $F_{\vec{\pmb{\b}}}$ and contributes $\tau$ under the re-scaling $k_{p_{m+1}}\to\tau k_{p_{m+1}}$. The combinatory momentum $k_{\{p_i\}\setminus\pmb{\a}}$ is the summation of momenta for external legs in the set $\{p_1,\cdots,p_m\}\setminus\pmb{\a}$. Applying the above procedure iteratively, one can expand the leading terms of YMS amplitudes in the second line of \eref{toB5-1} further, until there is only one gluon $p_{m+1}$ left. In other words, ${\cal A}^{(0)}_{\rm YS}(1,\{2,\cdots,n-1\}\shuffle\vec{\pmb{\a}},n;\{\{p_1,\cdots,p_m\}\setminus\pmb{\a}\}\cup p_{m+1}|\sigma_{n+m+1})$ can be expanded to
${\cal A}^{(0)}_{\rm YS}(1,\{2,\cdots,n-1\}\shuffle p_1\shuffle\cdots\shuffle p_m,n;p_{m+1}|\sigma_{n+m+1})$. Then, we expand them further by employing
%
\bea
& &{\cal A}^{(0)}_{\rm YS}(1,\{2,\cdots,n-1\}\shuffle p_1\shuffle\cdots\shuffle p_m,n;p_{m+1}|\sigma_{n+m+1})\nn
&=&{k_n\cdot f_{p_{m+1}}\cdot X_{p_{m+1}}\over k_n\cdot k_{p_{m+1}}}\,{\cal A}^{(0)}_{\rm S}(1,\{2,\cdots,n-1\}\shuffle p_1\shuffle\cdots\shuffle p_m\shuffle p_{m+1},n|\sigma_{n+m+1})\,,~~\label{toB5-2}
\eea
%
and reorganize the full expansion as
%
\bea
& &{\cal A}^{(0)}_{\rm YS}(1,\{2,\cdots,n-1\}\shuffle\vec{\pmb{\a}},n;\{\{p_1,\cdots,p_m\}\setminus\pmb{\a}\}\cup p_{m+1}|\sigma_{n+m+1})\nn
&=&{k_n\cdot f_{p_{m+1}}\cdot X_{p_{m+1}}\over k_n\cdot k_{p_{m+1}}}\,\W{\cal A}^{(0)}_{\rm YS}(1,\{2,\cdots,n-1\}\shuffle\vec{\pmb{\a}}\shuffle p_{m+1},n;\{p_1,\cdots,p_m\}\setminus\pmb{\a}|\sigma_{n+m+1})\,,~~\label{toB5-3}
\eea
%
based on the observation that $k_{p_{m+1}}$ which is accompanied with $\tau$ does not contribute to any $Y$ at the leading order.
The relation \eref{toB5-3} is the generalization of the previous result \eref{B83}.
Using \eref{toB5-3}, we arrive at the expression of $B_5$ in \eref{mg-B5}.

Now we observe that $B_1$ in \eref{mg-B1} together with $B_5$ in \eref{mg-B5} give the sub-leading term of
%
\bea
P_1=\sum_{\vec{\pmb{\a}}}\,{k_n\cdot F_{\vec{\pmb{\a}}}\cdot Y_{\vec{\pmb{\a}}}\over k_n\cdot k_{p_1\cdots p_{m+1}}}\,
{\cal A}_{\rm YS}(1,\{2,\cdots,n-1\}\shuffle\vec{\pmb{\a}},n;\{\{p_1,\cdots,p_m\}\setminus\pmb{\a}\}\cup p_{m+1}|\sigma_{n+m+1})\,,~~\label{mg-P1}
\eea
%
$B_2$ in \eref{mg-B2} can be interpreted as the leading term of
%
\bea
P_2=\sum_{\vec{\pmb{\a}}'}\,{k_n\cdot F_{\vec{\pmb{\a}}'}\cdot Y_{\vec{\pmb{\a}}'}\over k_n\cdot k_{p_1\cdots p_{m+1}}}\,
{\cal A}_{\rm YS}(1,\{2,\cdots,n-1\}\shuffle\vec{\pmb{\a}}',n;\{p_1,\cdots,p_m,p_{m+1}\}\setminus\pmb{\a'}|\sigma_{n+m+1})\,,~~\label{mg-P2}
\eea
%
where $\pmb{\a}'$ includes the external leg $p_{m+1}$ and $\pmb{\a}'\setminus p_{m+1}\neq\emptyset$, and $B_4$ in \eref{mg-B4} corresponds to the leading terms of
%
\bea
P_3={k_n\cdot f_{p_{m+1}}\cdot Y_{p_{m+1}}\over k_n\cdot k_{p_1\cdots p_{m+1}}}\,{\cal A}_{\rm YS}(1,\{2,\cdots,n-1\}\shuffle p_{m+1},n;p_1,\cdots,p_m|\sigma_{n+m+1})\,.~~\label{mg-P3}
\eea
%
Combining $P_1$, $P_2$ and $P_3$ together, we see that the general formula \eref{expan-YMS-gi} is correct for single trace YMS amplitude with $(m+1)$
external gluons, thus the proof is completed.


%%%%%%%%%%%%%%%%%%%%%%%%%%%%%
\section{Summery}
\label{sec-conclusion}
%%%%%%%%%%%%%%%%%%%%%%%%%%%%

In this note, we reproduced the expansion of single trace YMS amplitudes proposed by Clifford Cheung and James Mangan, via the soft bootstrap method. The whole process is on-shell, without using any off-shell ansatz. The resulting expansion manifests the gauge invariance for all polarizations carried by external gluons, with the cost of sacrificing the manifest locality. One can use such expansion iteratively to expand the single trace YMS amplitudes to pure BAS ones without breaking the manifest gauge invariance for any polarization, but a verity of spurious poles will be introduced. Via the double copy structure, our result can be extended to the expansion of single trace Einstein-Yang-Mills amplitudes directly,
%
\bea
& &{\cal A}_{\rm EY}(1,\cdots,n;p_1,\cdots,p_m)\nn
&=&\sum_{\vec{\pmb{\a}}}\,{k_n\cdot F_{\vec{\pmb{\a}}}\cdot Y_{\vec{\pmb{\a}}}\over k_n\cdot k_{p_1\cdots p_m}}\,{\cal A}_{\rm EY}(1,\{2,\cdots,n-1\}\shuffle\vec{\pmb{\a}},n;\{p_1,\cdots,p_m\}\setminus{\pmb{\a}})\,,~~\label{expan-EY-gi}
\eea
%
where $1,\cdots,n$ are color ordered gluons, and $p_1,\cdots,p_m$ are gravitons without any ordering. It is straightforward to verify
the expansion \eref{expan-EY-gi} is equivalent to that found by Clifford Cheung and James Mangan in \cite{Cheung:2021zvb}.

The sub-leading soft theorem for external gluon plays the crucial role in our method. One may ask the reason for choosing sub-leading soft theorem rather than the leading order one. The answer is, the leading order soft theorem can not detect all terms in the expansion. Suppose we re-scale $k_{p_j}$ as $k_{p_j}\to\tau k_{p_j}$, as can be seen in the general expanded formula \eref{expan-YMS-gi}, various terms include $f_{p_j}^{\mu\nu}$ in corresponding coefficients therefore are accompanied with $\tau$, these terms have no contribution to the leading order.
Another question is whether the sub-leading soft theorem is sufficient to detect all terms in the complete expansion? For example, if the coefficient of one term behaviors as $\tau^2$, then this coefficient can not be detected at both leading and sub-leading orders. This question was answered in \cite{Zhou:2022orv}, by employing the universality of soft factor for external scalar, as well as the counting of mass dimensions. The argument in \cite{Zhou:2022orv} at least shows that the old expansion \eref{expan-YMS} can be fully detected by the sub-leading soft theorem for external gluon. Since any correct expansion must be equivalent to \eref{expan-YMS}, we claim that the new expansion, which manifests gauge invariance for all polarizations, can also be detected completely.

A related interesting future direction is seeking the expanded formula for pure YM amplitudes which manifests the gauge invariance for each polarization. Such expansion leads to the manifestly gauge invariant BCJ numerators. Until now, such expansion has not been provided in any literature. We hope the method used in this note can be applied to solve this challenge.

\section*{Acknowledgments}

The authors would thank Prof.Yijian Du for stimulating discussions. We also thank Dr.Linghui Hou who told us the expansion found by Clifford Cheung and James Mangan.

%%%%%%%%%%%%%%%%%%%%%%%%%%%%%%%%%%%%%%%%%%%


\begin{thebibliography}{}
\bibliographystyle{JHEP}

\bibitem{Stieberger:2016lng}
  S.~Stieberger and T.~R.~Taylor,
  ``New relations for Einstein-Yang-Mills amplitudes,''
  Nucl.\ Phys.\ B {\bf 913}, 151 (2016)
 % doi:10.1016/j.nuclphysb.2016.09.014
  [arXiv:1606.09616 [hep-th]].
  %%CITATION = doi:10.1016/j.nuclphysb.2016.09.014;%%
  %39 citations counted in INSPIRE as of 18 Mar 2019


%\cite{Schlotterer:2016cxa}
\bibitem{Schlotterer:2016cxa}
  O.~Schlotterer,
  ``Amplitude relations in heterotic string theory and Einstein-Yang-Mills,''
  JHEP {\bf 1611}, 074 (2016)
  %doi:10.1007/JHEP11(2016)074
  [arXiv:1608.00130 [hep-th]].
  %%CITATION = doi:10.1007/JHEP11(2016)074;%%
  %23 citations counted in INSPIRE as of 18 Mar 2019


%\cite{Chiodaroli:2017ngp}
\bibitem{Chiodaroli:2017ngp}
  M.~Chiodaroli, M.~Gunaydin, H.~Johansson and R.~Roiban,
  ``Explicit Formulae for Yang-Mills-Einstein Amplitudes from the Double Copy,''
  JHEP {\bf 1707}, 002 (2017)
  doi:10.1007/JHEP07(2017)002
  [arXiv:1703.00421 [hep-th]].
  %%CITATION = doi:10.1007/JHEP07(2017)002;%%
  %35 citations counted in INSPIRE as of 31 May 2019

%\cite{DelDuca:1999rs}
\bibitem{DelDuca:1999rs}
  V.~Del Duca, L.~J.~Dixon and F.~Maltoni,
  ``New color decompositions for gauge amplitudes at tree and loop level,''
  Nucl.\ Phys.\ B {\bf 571}, 51 (2000)
  doi:10.1016/S0550-3213(99)00809-3
  [hep-ph/9910563].
  %%CITATION = doi:10.1016/S0550-3213(99)00809-3;%%
  %225 citations counted in INSPIRE as of 06 Jun 2019


%\cite{Nandan:2016pya}
\bibitem{Nandan:2016pya}
  D.~Nandan, J.~Plefka, O.~Schlotterer and C.~Wen,
  ``Einstein-Yang-Mills from pure Yang-Mills amplitudes,''
  JHEP {\bf 1610}, 070 (2016)
  %doi:10.1007/JHEP10(2016)070
  [arXiv:1607.05701 [hep-th]].
  %%CITATION = doi:10.1007/JHEP10(2016)070;%%
  %34 citations counted in INSPIRE as of 18 Mar 2019


%\cite{delaCruz:2016gnm}
\bibitem{delaCruz:2016gnm}
  L.~de la Cruz, A.~Kniss and S.~Weinzierl,
  ``Relations for Einstein-Yang-Mills amplitudes from the CHY representation,''
  Phys.\ Lett.\ B {\bf 767}, 86 (2017)
 % doi:10.1016/j.physletb.2017.01.036
  [arXiv:1607.06036 [hep-th]].
  %%CITATION = doi:10.1016/j.physletb.2017.01.036;%%
  %25 citations counted in INSPIRE as of 18 Mar 2019

%\cite{Fu:2017uzt}
\bibitem{Fu:2017uzt}
  C.~H.~Fu, Y.~J.~Du, R.~Huang and B.~Feng,
  ``Expansion of Einstein-Yang-Mills Amplitude,''
  JHEP {\bf 1709}, 021 (2017)
  doi:10.1007/JHEP09(2017)021
  [arXiv:1702.08158 [hep-th]].
  %%CITATION = doi:10.1007/JHEP09(2017)021;%%
  %11 citations counted in INSPIRE as of 12 Aug 2018


%\cite{Fu:2017uzt,Teng:2017tbo,Du:2017kpo,Du:2017gnh,Feng:2019tvb,Zhou:2022orv}
\bibitem{Teng:2017tbo}
  F.~Teng and B.~Feng,
  ``Expanding Einstein-Yang-Mills by Yang-Mills in CHY frame,''
  JHEP {\bf 1705}, 075 (2017)
  %doi:10.1007/JHEP05(2017)075
  [arXiv:1703.01269 [hep-th]].
  %%CITATION = doi:10.1007/JHEP05(2017)075;%%
  %19 citations counted in INSPIRE as of 18 Mar 2019

%\cite{Du:2017kpo,Du:2017gnh,Feng:2019tvb}
\bibitem{Du:2017kpo}
  Y.~J.~Du and F.~Teng,
  ``BCJ numerators from reduced Pfaffian,''
  JHEP {\bf 1704}, 033 (2017)
 % doi:10.1007/JHEP04(2017)033
  [arXiv:1703.05717 [hep-th]].
  %%CITATION = doi:10.1007/JHEP04(2017)033;%%
  %19 citations counted in INSPIRE as of 18 Mar 2019


%\cite{Du:2017gnh,Feng:2019tvb}
\bibitem{Du:2017gnh}
  Y.~J.~Du, B.~Feng and F.~Teng,
  ``Expansion of All Multitrace Tree Level EYM Amplitudes,''
  JHEP {\bf 1712}, 038 (2017)
 % doi:10.1007/JHEP12(2017)038
  [arXiv:1708.04514 [hep-th]].
  %%CITATION = doi:10.1007/JHEP12(2017)038;%%
  %16 citations counted in INSPIRE as of 18 Mar 2019

%\cite{Feng:2019tvb}
\bibitem{Feng:2019tvb}
  B.~Feng, X.~Li and K.~Zhou,
  ``Expansion of EYM theory by Differential Operators,''
  arXiv:1904.05997 [hep-th].
  %%CITATION = ARXIV:1904.05997;%%

%\cite{Zhou:2022orv}
\bibitem{Zhou:2022orv}
K.~Zhou,
``Tree level amplitudes from soft theorems,''
JHEP \textbf{03}, 021 (2023)
doi:10.1007/JHEP03(2023)021
[arXiv:2212.12892 [hep-th]].
%0 citations counted in INSPIRE as of 01 May 2023

%%%%%%%%%%%%%%%%%%%%%%%%%%%%%%%%%%%%%%%%%%%%%%%%%%%%%%%%

%\cite{Cheung:2021zvb}
\bibitem{Cheung:2021zvb}
C.~Cheung and J.~Mangan,
``Covariant color-kinematics duality,''
JHEP \textbf{11}, 069 (2021)
doi:10.1007/JHEP11(2021)069
[arXiv:2108.02276 [hep-th]].
%44 citations counted in INSPIRE as of 01 Jun 2023

%%%%%%%%%%%%%%%%%%%%%%%%%%%%%%%%%%%%%%%%%%%%%%%%%%%%%%%%%

%\cite{Bern:2008qj,Chiodaroli:2014xia,Johansson:2015oia,Johansson:2019dnu}
\bibitem{Bern:2008qj}
  Z.~Bern, J.~J.~M.~Carrasco and H.~Johansson,
  ``New Relations for Gauge-Theory Amplitudes,''
  Phys.\ Rev.\ D {\bf 78}, 085011 (2008)
  %doi:10.1103/PhysRevD.78.085011
  [arXiv:0805.3993 [hep-ph]].
  %%CITATION = doi:10.1103/PhysRevD.78.085011;%%
  %578 citations counted in INSPIRE as of 18 Mar 2019

%\cite{Chiodaroli:2014xia,Johansson:2015oia,Johansson:2019dnu}
\bibitem{Chiodaroli:2014xia}
  M.~Chiodaroli, M.~Gnaydin, H.~Johansson and R.~Roiban,
  ``Scattering amplitudes in $ \mathcal{N}=2 $ Maxwell-Einstein and Yang-Mills/Einstein supergravity,''
  JHEP {\bf 1501}, 081 (2015)
  doi:10.1007/JHEP01(2015)081
  [arXiv:1408.0764 [hep-th]].

%\cite{Johansson:2015oia,Johansson:2019dnu}
\bibitem{Johansson:2015oia}
  H.~Johansson and A.~Ochirov,
  ``Color-Kinematics Duality for QCD Amplitudes,''
  JHEP {\bf 1601}, 170 (2016)
  doi:10.1007/JHEP01(2016)170
  [arXiv:1507.00332 [hep-ph]].

%\cite{Johansson:2019dnu}
\bibitem{Johansson:2019dnu}
  H.~Johansson and A.~Ochirov,
  ``Double copy for massive quantum particles with spin,''
  JHEP {\bf 1909}, 040 (2019)
  doi:10.1007/JHEP09(2019)040
  [arXiv:1906.12292 [hep-th]].
  %%CITATION = doi:10.1007/JHEP09(2019)040;%%

%%%%%%%%%%%%%%%%%%%%%%%%%%%%%%%%%%%%%%%%%%%%%%%%%%%%%%

%\cite{Casali:2014xpa,Schwab:2014xua,,Afkhami-Jeddi:2014fia}
\bibitem{Casali:2014xpa}
E.~Casali,
``Soft sub-leading divergences in Yang-Mills amplitudes,''
JHEP \textbf{08}, 077 (2014)
doi:10.1007/JHEP08(2014)077
[arXiv:1404.5551 [hep-th]].
%138 citations counted in INSPIRE as of 15 Dec 2022

%\cite{Schwab:2014xua}
\bibitem{Schwab:2014xua}
B.~U.~W.~Schwab and A.~Volovich,
``Subleading Soft Theorem in Arbitrary Dimensions from Scattering Equations,''
Phys. Rev. Lett. \textbf{113}, no.10, 101601 (2014)
doi:10.1103/PhysRevLett.113.101601
[arXiv:1404.7749 [hep-th]].
%139 citations counted in INSPIRE as of 15 Dec 2022

%\cite{Afkhami-Jeddi:2014fia}
\bibitem{Afkhami-Jeddi:2014fia}
N.~Afkhami-Jeddi,
``Soft Graviton Theorem in Arbitrary Dimensions,''
[arXiv:1405.3533 [hep-th]].
%98 citations counted in INSPIRE as of 15 Dec 2022

%%%%%%%%%%%%%%%%%%%%%%%%%%%%%%%%%%%%%%%%%%%%%%%%%%%%%

\bibitem{Cachazo:2013iea}
F.~Cachazo, S.~He, and E.~Y. Yuan,
``{Scattering of Massless Particles: Scalars, Gluons and Gravitons},''
  { JHEP} {\bf 1407} (2014)
   033,
 arXiv:1309.0885 [hep-th].
%%CITATION = ARXIV:1309.0885;%%.

%%%%%%%%%%%%%%%%%%%%%%%%%%%%%%%%%%%%%%%%%%%%%%%%%%

\bibitem{Kawai:1985xq}
  H.~Kawai, D.~C.~Lewellen and S.~H.~Tye,
  ``A Relation Between Tree Amplitudes of Closed and Open Strings,''
  Nucl.\ Phys.\ B {\bf 269}, 1 (1986).
 % doi:10.1016/0550-3213(86)90362-7
  %%CITATION = doi:10.1016/0550-3213(86)90362-7;%%
  %697 citations counted in INSPIRE as of 31 May 2019

%%%%%%%%%%%%%%%%%%%%%%%%%%%%%%%%%%%%%%%%%%%%%%%%%%%%%

%\cite{Kleiss:1988ne}
\bibitem{Kleiss:1988ne}
  R.~Kleiss and H.~Kuijf,
  ``MULTI - GLUON CROSS-SECTIONS AND FIVE JET PRODUCTION AT HADRON COLLIDERS,''
  Nucl.\ Phys.\  B {\bf 312}, 616 (1989).
  %%CITATION = NUPHA,B312,616;%%

%%%%%%%%%%%%%%%%%%%%%%%%%%%%%%%%%%%%%%%%%%

\end{thebibliography}


\end{document}
